% Options for packages loaded elsewhere
\PassOptionsToPackage{unicode}{hyperref}
\PassOptionsToPackage{hyphens}{url}
\documentclass[
  11pt,
]{article}
\usepackage{xcolor}
\usepackage[left=3cm,right=5cm,top=2cm,bottom=2cm]{geometry}
\usepackage{amsmath,amssymb}
\setcounter{secnumdepth}{5}
\usepackage{iftex}
\ifPDFTeX
  \usepackage[T1]{fontenc}
  \usepackage[utf8]{inputenc}
  \usepackage{textcomp} % provide euro and other symbols
\else % if luatex or xetex
  \usepackage{unicode-math} % this also loads fontspec
  \defaultfontfeatures{Scale=MatchLowercase}
  \defaultfontfeatures[\rmfamily]{Ligatures=TeX,Scale=1}
\fi
\usepackage{lmodern}
\ifPDFTeX\else
  % xetex/luatex font selection
\fi
% Use upquote if available, for straight quotes in verbatim environments
\IfFileExists{upquote.sty}{\usepackage{upquote}}{}
\IfFileExists{microtype.sty}{% use microtype if available
  \usepackage[]{microtype}
  \UseMicrotypeSet[protrusion]{basicmath} % disable protrusion for tt fonts
}{}
\makeatletter
\@ifundefined{KOMAClassName}{% if non-KOMA class
  \IfFileExists{parskip.sty}{%
    \usepackage{parskip}
  }{% else
    \setlength{\parindent}{0pt}
    \setlength{\parskip}{6pt plus 2pt minus 1pt}}
}{% if KOMA class
  \KOMAoptions{parskip=half}}
\makeatother
\usepackage{color}
\usepackage{fancyvrb}
\newcommand{\VerbBar}{|}
\newcommand{\VERB}{\Verb[commandchars=\\\{\}]}
\DefineVerbatimEnvironment{Highlighting}{Verbatim}{commandchars=\\\{\}}
% Add ',fontsize=\small' for more characters per line
\usepackage{framed}
\definecolor{shadecolor}{RGB}{248,248,248}
\newenvironment{Shaded}{\begin{snugshade}}{\end{snugshade}}
\newcommand{\AlertTok}[1]{\textcolor[rgb]{0.94,0.16,0.16}{#1}}
\newcommand{\AnnotationTok}[1]{\textcolor[rgb]{0.56,0.35,0.01}{\textbf{\textit{#1}}}}
\newcommand{\AttributeTok}[1]{\textcolor[rgb]{0.13,0.29,0.53}{#1}}
\newcommand{\BaseNTok}[1]{\textcolor[rgb]{0.00,0.00,0.81}{#1}}
\newcommand{\BuiltInTok}[1]{#1}
\newcommand{\CharTok}[1]{\textcolor[rgb]{0.31,0.60,0.02}{#1}}
\newcommand{\CommentTok}[1]{\textcolor[rgb]{0.56,0.35,0.01}{\textit{#1}}}
\newcommand{\CommentVarTok}[1]{\textcolor[rgb]{0.56,0.35,0.01}{\textbf{\textit{#1}}}}
\newcommand{\ConstantTok}[1]{\textcolor[rgb]{0.56,0.35,0.01}{#1}}
\newcommand{\ControlFlowTok}[1]{\textcolor[rgb]{0.13,0.29,0.53}{\textbf{#1}}}
\newcommand{\DataTypeTok}[1]{\textcolor[rgb]{0.13,0.29,0.53}{#1}}
\newcommand{\DecValTok}[1]{\textcolor[rgb]{0.00,0.00,0.81}{#1}}
\newcommand{\DocumentationTok}[1]{\textcolor[rgb]{0.56,0.35,0.01}{\textbf{\textit{#1}}}}
\newcommand{\ErrorTok}[1]{\textcolor[rgb]{0.64,0.00,0.00}{\textbf{#1}}}
\newcommand{\ExtensionTok}[1]{#1}
\newcommand{\FloatTok}[1]{\textcolor[rgb]{0.00,0.00,0.81}{#1}}
\newcommand{\FunctionTok}[1]{\textcolor[rgb]{0.13,0.29,0.53}{\textbf{#1}}}
\newcommand{\ImportTok}[1]{#1}
\newcommand{\InformationTok}[1]{\textcolor[rgb]{0.56,0.35,0.01}{\textbf{\textit{#1}}}}
\newcommand{\KeywordTok}[1]{\textcolor[rgb]{0.13,0.29,0.53}{\textbf{#1}}}
\newcommand{\NormalTok}[1]{#1}
\newcommand{\OperatorTok}[1]{\textcolor[rgb]{0.81,0.36,0.00}{\textbf{#1}}}
\newcommand{\OtherTok}[1]{\textcolor[rgb]{0.56,0.35,0.01}{#1}}
\newcommand{\PreprocessorTok}[1]{\textcolor[rgb]{0.56,0.35,0.01}{\textit{#1}}}
\newcommand{\RegionMarkerTok}[1]{#1}
\newcommand{\SpecialCharTok}[1]{\textcolor[rgb]{0.81,0.36,0.00}{\textbf{#1}}}
\newcommand{\SpecialStringTok}[1]{\textcolor[rgb]{0.31,0.60,0.02}{#1}}
\newcommand{\StringTok}[1]{\textcolor[rgb]{0.31,0.60,0.02}{#1}}
\newcommand{\VariableTok}[1]{\textcolor[rgb]{0.00,0.00,0.00}{#1}}
\newcommand{\VerbatimStringTok}[1]{\textcolor[rgb]{0.31,0.60,0.02}{#1}}
\newcommand{\WarningTok}[1]{\textcolor[rgb]{0.56,0.35,0.01}{\textbf{\textit{#1}}}}
\usepackage{longtable,booktabs,array}
\newcounter{none} % for unnumbered tables
\usepackage{calc} % for calculating minipage widths
% Correct order of tables after \paragraph or \subparagraph
\usepackage{etoolbox}
\makeatletter
\patchcmd\longtable{\par}{\if@noskipsec\mbox{}\fi\par}{}{}
\makeatother
% Allow footnotes in longtable head/foot
\IfFileExists{footnotehyper.sty}{\usepackage{footnotehyper}}{\usepackage{footnote}}
\makesavenoteenv{longtable}
\usepackage{graphicx}
\makeatletter
\newsavebox\pandoc@box
\newcommand*\pandocbounded[1]{% scales image to fit in text height/width
  \sbox\pandoc@box{#1}%
  \Gscale@div\@tempa{\textheight}{\dimexpr\ht\pandoc@box+\dp\pandoc@box\relax}%
  \Gscale@div\@tempb{\linewidth}{\wd\pandoc@box}%
  \ifdim\@tempb\p@<\@tempa\p@\let\@tempa\@tempb\fi% select the smaller of both
  \ifdim\@tempa\p@<\p@\scalebox{\@tempa}{\usebox\pandoc@box}%
  \else\usebox{\pandoc@box}%
  \fi%
}
% Set default figure placement to htbp
\def\fps@figure{htbp}
\makeatother
% definitions for citeproc citations
\NewDocumentCommand\citeproctext{}{}
\NewDocumentCommand\citeproc{mm}{%
  \begingroup\def\citeproctext{#2}\cite{#1}\endgroup}
\makeatletter
 % allow citations to break across lines
 \let\@cite@ofmt\@firstofone
 % avoid brackets around text for \cite:
 \def\@biblabel#1{}
 \def\@cite#1#2{{#1\if@tempswa , #2\fi}}
\makeatother
\newlength{\cslhangindent}
\setlength{\cslhangindent}{1.5em}
\newlength{\csllabelwidth}
\setlength{\csllabelwidth}{3em}
\newenvironment{CSLReferences}[2] % #1 hanging-indent, #2 entry-spacing
 {\begin{list}{}{%
  \setlength{\itemindent}{0pt}
  \setlength{\leftmargin}{0pt}
  \setlength{\parsep}{0pt}
  % turn on hanging indent if param 1 is 1
  \ifodd #1
   \setlength{\leftmargin}{\cslhangindent}
   \setlength{\itemindent}{-1\cslhangindent}
  \fi
  % set entry spacing
  \setlength{\itemsep}{#2\baselineskip}}}
 {\end{list}}
\usepackage{calc}
\newcommand{\CSLBlock}[1]{\hfill\break\parbox[t]{\linewidth}{\strut\ignorespaces#1\strut}}
\newcommand{\CSLLeftMargin}[1]{\parbox[t]{\csllabelwidth}{\strut#1\strut}}
\newcommand{\CSLRightInline}[1]{\parbox[t]{\linewidth - \csllabelwidth}{\strut#1\strut}}
\newcommand{\CSLIndent}[1]{\hspace{\cslhangindent}#1}
\setlength{\emergencystretch}{3em} % prevent overfull lines
\providecommand{\tightlist}{%
  \setlength{\itemsep}{0pt}\setlength{\parskip}{0pt}}
% bullets  -- boiled-down topic sentence + bullet points
%             small, sans-serif, dark gray; paragraph number [N] prepended
% rgt      -- author's rewrite (initially a placeholder)
%             normal size, dark blue, serif
% orig     -- original Claude-authored draft
%             small, italic, light gray

\usepackage{xcolor}

\newcounter{pgraph}

\newenvironment{bullets}{%
  \stepcounter{pgraph}%
  \begin{quote}\small\sffamily\color[gray]{0.30}%
  {\normalfont\scriptsize\color[gray]{0.58}[\thepgraph]\enspace}}{%
  \end{quote}}

\newenvironment{rgt}{%
  \begin{quote}\normalsize\color[RGB]{20,60,160}}{%
  \end{quote}}

\newenvironment{orig}{%
  \begin{quote}\small\itshape\color[gray]{0.55}}{%
  \end{quote}}

% backward compat alias
\newenvironment{claudecode}{\begin{orig}}{\end{orig}}
\usepackage{amsmath}
\usepackage{amssymb}
\usepackage{lineno}
\linenumbers
\usepackage{booktabs}
\usepackage{float}
\DeclareMathOperator{\Var}{Var}
\DeclareMathOperator{\Cov}{Cov}
\DeclareMathOperator{\tr}{tr}
\usepackage{setspace}
\setstretch{1.1}
\usepackage{bookmark}
\IfFileExists{xurl.sty}{\usepackage{xurl}}{} % add URL line breaks if available
\urlstyle{same}
% fallback for those not using the hyperref driver hyperxmp:
\makeatletter
\@ifundefined{xmpquote}{\newcommand{\xmpquote}[1]{#1}}{}
\makeatother
\hypersetup{
  pdftitle={Power Analysis for Longitudinal Clinical Trials with Run-In and Common Close Observations},
  pdfauthor={Ronald G. Thomas, Ph.D.},
  hidelinks,
  pdfcreator={LaTeX via pandoc}}

\title{Power Analysis for Longitudinal Clinical Trials
with Run-In and Common Close Observations}
\author{Ronald G. Thomas, Ph.D.\thanks{Wertheim School of Public Health, UCSD; ORCID: 0000-0003-1686-4965}}
\date{July 01, 2026}

\begin{document}
\maketitle
\begin{abstract}
\textbf{Background.} Longitudinal clinical trials often
incorporate pre-randomisation (run-in) and post-treatment
(common close) observations to improve estimator efficiency.
Closed-form power and sample-size formulas for these
augmented designs under linear mixed-effects models with
random slopes have remained scattered across the literature
and partially specified, making applied sample-size
calculation cumbersome.

\textbf{Methods.} We derive closed-form power and sample-size
formulas for three nested cases: (i) two run-in observations,
(ii) an arbitrary number \(J_0\) of run-in observations, and
(iii) \(J_0\) run-in plus \(J_2\) common close observations. All
derivations proceed via the Woodbury matrix identity applied
to the GLS estimator under change-score parameterisation.
The formulas are expressed in terms of the residual variance
\(\sigma^2\), the random-slope variance \(\sigma_b^2\), and the
observation spacing \(t\), enabling direct computation of
required sample sizes.

\textbf{Results.} In a numerical example using parameters from
Alzheimer's disease neuroimaging trials, for which the
signal-to-noise ratio \(\sigma_b / \sigma \approx 1.7\) is
favourable, two run-in observations reduce the required
sample size by 62\% (from 385 to 145 per group), and adding
two common close observations raises the reduction to 82\%.
The magnitude of the gain is governed by \(\sigma_b / \sigma\):
across a representative range of slope variability the
reduction from two run-in observations spans roughly 14\% at
low ratios to 76\% at high ratios. The closed-form
expressions allow this sensitivity analysis with respect to
\(\sigma_b / \sigma\) without re-running simulation.

\textbf{Conclusions.} The Woodbury-based formulas are the simplest
available closed-form expressions for the variance of the
treatment-effect estimator under random-slopes
longitudinal-trial designs with run-in and common close
observations. They serve as a foundation for the AR(1)
extension (Paper 2), the optimal-allocation problem
(Paper 3), and the GLS-versus-ANCOVA-versus-averaged
estimator comparison (Paper 4) in this compendium.
\end{abstract}

{
\setcounter{tocdepth}{2}
\tableofcontents
}
\section{Introduction}\label{introduction}

\begin{bullets}
\textbf{Beyond its traditional roles, the run-in period has a distinct methodological value in rate-of-change trials: estimating individual baseline trajectories to raise power.}
\begin{itemize}\setlength{\itemsep}{2pt}\setlength{\parskip}{0pt}
\item Conventional run-in uses include identifying compliant participants and washing out prior treatments.
\item When the outcome is a rate of change, run-in observations estimate each participant's baseline trajectory.
\item Incorporating those trajectories into the analysis model substantially increases statistical power.
\end{itemize}
\end{bullets}

\begin{rgt}
To be completed by rgt.
\end{rgt}

\begin{orig}
The run-in period has served multiple purposes in clinical trial
design, from identifying compliant participants to washing out
prior treatments\textsuperscript{\citeproc{ref-packer2017}{1}}. A distinct and
methodologically important use has emerged in trials where the
primary outcome is a rate of change: estimating individual
baseline trajectories that, when incorporated into the analysis
model, substantially increase statistical power.

\end{orig}

\begin{bullets}
\textbf{Frost and colleagues established the mixed-model foundation: run-in observations help most when slope variability is large relative to measurement error.}
\begin{itemize}\setlength{\itemsep}{2pt}\setlength{\parskip}{0pt}
\item The key insight is that pre-randomization observations inform each participant's underlying trajectory.
\item This reduces the residual variance of the treatment-effect estimator, and the gain grows with the slope-to-noise ratio.
\item In Alzheimer's neuroimaging trials, where that ratio is favorable, they reported sample-size reductions of 30 to 45 percent.
\end{itemize}
\end{bullets}

\begin{rgt}
To be completed by rgt.
\end{rgt}

\begin{orig}
\textsuperscript{\citeproc{ref-frost2008}{2}} established the theoretical
foundation for this approach within the linear mixed effects
framework. Their key insight was that when between-subject
variability in rates of change is large relative to measurement
error, pre-randomization observations provide valuable
information about each participant's underlying trajectory,
reducing the residual variance of the treatment effect
estimator. Working with neuroimaging outcomes in Alzheimer's
disease trials, where the ratio of slope variability to
measurement error is particularly favorable, they demonstrated
sample size reductions of 30--45\%.

\end{orig}

\begin{bullets}
\textbf{Wang and colleagues extended the framework with two-period models that jointly model run-in and post-randomization data rather than treating run-in as a covariate.}
\begin{itemize}\setlength{\itemsep}{2pt}\setlength{\parskip}{0pt}
\item Joint modelling of both periods yields additional power gains of up to 15 percent.
\item It permits unequal randomization ratios without power loss.
\item Run-in observations contribute placebo-equivalent information for all participants.
\end{itemize}
\end{bullets}

\begin{rgt}
To be completed by rgt.
\end{rgt}

\begin{orig}
\textsuperscript{\citeproc{ref-wang2019}{3}} extended this framework by proposing
two-period linear mixed effects models that jointly model
run-in and post-randomization data rather than using run-in
observations merely as covariates. Their approach yields
additional power gains of up to 15\% and enables unequal
randomization ratios without power loss, since run-in
observations effectively contribute placebo-equivalent
information for all participants.

\end{orig}

\begin{bullets}
\textbf{A parallel literature gives closed-form variance and sample-size formulae for the random-coefficient model without a run-in phase, which form the baseline this paper extends.}
\begin{itemize}\setlength{\itemsep}{2pt}\setlength{\parskip}{0pt}
\item Lu and colleagues derived the cLDA estimator variance, its link to ANCOVA, and the conditions under which the two coincide.
\item Hu and colleagues generalized these to random-coefficient models with monotone missing data and variable follow-up, illustrated with ADNI data.
\item These results cover only the $J_0 = 0$ case and serve as the baseline extended here.
\end{itemize}
\end{bullets}

\begin{rgt}
To be completed by rgt.
\end{rgt}

\begin{orig}
A parallel literature has developed closed-form variance and
sample size formulae for the random coefficient regression
model without a run-in phase.\textsuperscript{\citeproc{ref-lu2009}{4}} derived the variance of
the constrained longitudinal data analysis estimator and its
relation to analysis of covariance, and\textsuperscript{\citeproc{ref-lu2010}{5}} established
efficiency conditions under which the two approaches coincide.\textsuperscript{\citeproc{ref-hu2021}{6}} generalized these results to random coefficient
regression models with monotone missing data and variable
follow-up, providing variance, power, and sample size
expressions illustrated with data from the Alzheimer's Disease
Neuroimaging Initiative. These formulae apply to designs with
\(J_0 = 0\) run-in observations and an arbitrary post-randomization
visit schedule, and serve as the baseline case that we extend here.

\end{orig}

\begin{bullets}
\textbf{The CRAN longpower package is the reference implementation of the baseline formulae, and our general result reduces to its routines at the appropriate boundaries.}
\begin{itemize}\setlength{\itemsep}{2pt}\setlength{\parskip}{0pt}
\item longpower provides edland, diggle, and liu.liang routines for slope-comparison trials.
\item Our formula reduces to the Ard-Edland slope variance at $J_0 = J_2 = 0$, and to the Diggle compound-symmetry form when also $\sigma_b^2 = 0$.
\item The generalization thus embeds the existing longpower formulae as the baseline against which run-in and common-close gains are measured.
\end{itemize}
\end{bullets}

\begin{rgt}
To be completed by rgt.
\end{rgt}

\begin{orig}
The CRAN reference implementation of the \(J_0=0\) formulae is
the \texttt{longpower} package\textsuperscript{\citeproc{ref-iddi2022}{7}}, which provides the
routines \texttt{edland.linear.power()}, \texttt{diggle.linear.power()}, and
\texttt{liu.liang.linear.power()} for slope-comparison trials with
random intercepts and slopes\textsuperscript{\citeproc{ref-ard2011}{8}--\citeproc{ref-liuliang1997}{10}}.
We show in Section 3.3 that the closed-form result derived below
reduces to the Ard-Edland slope-variance formula implemented in
\texttt{edland.linear.power()} at the boundary \(J_0 = J_2 = 0\), and
to the Diggle compound-symmetry formula when additionally
\(\sigma_b^2 = 0\). The generalization developed here therefore
embeds the existing \texttt{longpower} formulae as the baseline case
against which the efficiency gains from run-in and common
close observations are measured.

\end{orig}

\begin{bullets}
\textbf{No prior closed-form power formula handles an arbitrary number of run-in observations together with a separate common-close phase, the gap this paper fills.}
\begin{itemize}\setlength{\itemsep}{2pt}\setlength{\parskip}{0pt}
\item The common-close phase observes all participants off-treatment after the randomized period.
\item It supplies extra information about individual trajectories at no cost to the treatment comparison.
\item Such post-treatment periods appear in delayed-start and randomized-withdrawal designs, but their GLS efficiency contribution has not been formally characterized.
\end{itemize}
\end{bullets}

\begin{rgt}
To be completed by rgt.
\end{rgt}

\begin{orig}
Despite these advances, no general closed-form power formula
has been derived that simultaneously accounts for an arbitrary
number of run-in observations and a separate set of common
close observations collected after treatment cessation. The
common close phase, in which all participants are observed
off-treatment following the randomized period, provides
additional information about individual trajectories at no
cost to the treatment comparison. Post-treatment observation
periods appear in delayed-start and randomized withdrawal
designs aimed at distinguishing symptomatic from
disease-modifying effects\textsuperscript{\citeproc{ref-sperling2014}{11}},
but the contribution of the common close to statistical
efficiency within the GLS framework has not been formally
characterized.

\end{orig}

\begin{bullets}
\textbf{The paper presents three nested results, all derived from the GLS estimator via the Woodbury identity.}
\begin{itemize}\setlength{\itemsep}{2pt}\setlength{\parskip}{0pt}
\item A closed-form variance for exactly two run-in observations (Section 3.2), then its generalization to arbitrary $J_0$ (Section 3.3).
\item A further extension adding $J_2$ common-close observations (Section 3.4).
\item All expressions are in terms of $\sigma^2$, $\sigma_b^2$, and spacing $t$, with numerical examples from Alzheimer's neuroimaging studies (Section 5).
\end{itemize}
\end{bullets}

\begin{rgt}
To be completed by rgt.
\end{rgt}

\begin{orig}
We present three results. First, we derive a simple
closed-form variance formula for the treatment effect
estimator when exactly two run-in observations precede
randomization (Section 3.2). Second, we generalize this
formula to an arbitrary number \(J_0\) of run-in observations
(Section 3.3). Third, we further extend to include \(J_2\) common
close observations collected after treatment cessation
(Section 3.4). All derivations proceed from the generalized
least squares estimator under the Woodbury matrix identity,
yielding expressions in terms of the residual variance
\(\sigma^2\), the random slope variance \(\sigma_b^2\), and the
observation spacing \(t\). Section 5 provides numerical
examples using parameters estimated from Alzheimer's disease
neuroimaging studies.

\end{orig}

\section{Model}\label{model}

\subsection{General formulation}\label{general-formulation}

\begin{bullets}
\textbf{The model accommodates a two-arm longitudinal trial whose observation schedule spans three distinct phases.}
\begin{itemize}\setlength{\itemsep}{2pt}\setlength{\parskip}{0pt}
\item $N$ participants are observed at equally spaced time points.
\item The phases are run-in (pre-randomization), treatment (post-randomization), and an optional common close (post-treatment).
\item $Y_{ij}$ denotes the outcome for participant $i$ at time point $j$.
\end{itemize}
\end{bullets}

\begin{rgt}
To be completed by rgt.
\end{rgt}

\begin{orig}
Consider a two-arm longitudinal clinical trial in which \(N\)
participants are observed at equally spaced time points across
three phases: a run-in period (pre-randomization), a treatment
period (post-randomization), and an optional common close
period (post-treatment). Let \(Y_{ij}\) denote the outcome for
participant \(i\) at time point \(j\). The observation schedule is:

\end{orig}

\begin{itemize}
\tightlist
\item
  \textbf{Run-in:} \(j = -J_0, -J_0+1, \ldots, -1\) (all participants
  untreated)
\item
  \textbf{Randomization visit:} \(j = 0\) (baseline)
\item
  \textbf{Treatment period:} \(j = 1, 2, \ldots, J_1\) (participants
  randomized to treatment or placebo)
\item
  \textbf{Common close:} \(j = J_1+1, J_1+2, \ldots, J_1+J_2\) (all
  participants off-treatment)
\end{itemize}

We model the outcome as
\begin{equation}
  Y_{ij} = \alpha + a_i +
  \bigl(\delta(1-h_j) + \beta h_j + \gamma h_j g_i + b_i
  \bigr) t_j + u_{ij} + \varepsilon_{ij}
  \label{eq:model}
\end{equation}

\begin{bullets}
\textbf{The outcome model combines a random intercept and random slope with a three-component within-subject variance structure following Frost and colleagues.}
\begin{itemize}\setlength{\itemsep}{2pt}\setlength{\parskip}{0pt}
\item $a_i$ and $b_i$ are the random intercept and slope, while $u_{ij}$ and $\varepsilon_{ij}$ are the between-visit and between-scan residuals.
\item Both residual terms are independent across visits, so under a single scan per visit they enter every formula only through their sum.
\item We therefore define the total within-visit residual $w_{ij} = u_{ij} + \varepsilon_{ij}$ with variance $\sigma^2$.
\end{itemize}
\end{bullets}

\begin{rgt}
To be completed by rgt.
\end{rgt}

\begin{orig}
where \(\alpha\) is the overall intercept, \(a_i \sim N(0,
\sigma_a^2)\) is a random intercept, \(b_i \sim N(0,
\sigma_b^2)\) is a random slope, \(u_{ij} \sim N(0,
\sigma_u^2)\) are within-subject between-visit random
effects independent across visits, and \(\varepsilon_{ij}
\sim N(0, \sigma_\varepsilon^2)\) are independent
between-scan measurement errors. This three-component
within-subject variance structure follows\textsuperscript{\citeproc{ref-frost2008}{2}} (their Eq. 17). Because
\(u_{ij}\) and \(\varepsilon_{ij}\) are both independent across
visits, under a single scan per visit they enter every
formula below only through their sum. We therefore write
\(w_{ij} = u_{ij} + \varepsilon_{ij}\) for the total
within-visit residual, with

\end{orig}

\begin{equation}
  \sigma^2 \;\equiv\; \mathop{\mathrm{Var}}(w_{ij}) \;=\; \sigma_u^2 +
  \sigma_\varepsilon^2 .
  \label{eq:resid}
\end{equation}

\begin{bullets}
\textbf{The two within-visit components are identifiable only with replicate scans, so single-scan designs depend solely on their sum.}
\begin{itemize}\setlength{\itemsep}{2pt}\setlength{\parskip}{0pt}
\item Separating $\sigma_u^2$ from $\sigma_\varepsilon^2$ requires a visit carrying replicate scans.
\item For the single-scan designs treated here, only the sum $\sigma^2$ enters the treatment-effect variance.
\item This reduction simplifies all subsequent variance expressions.
\end{itemize}
\end{bullets}

\begin{rgt}
To be completed by rgt.
\end{rgt}

\begin{orig}
The two within-visit components are separately identifiable
only when a visit carries replicate scans; for the
single-scan designs treated here only their sum \(\sigma^2\)
enters the variance of the treatment-effect estimator. The
phase indicator is

\end{orig}

\[
  h_j =
  \begin{cases}
    0 & \text{if } j \le 0 \quad \text{(run-in)} \\
    1 & \text{if } j > 0 \quad \text{(post-randomization)}
  \end{cases}
\]
and the treatment group indicator is
\[
  g_i =
  \begin{cases}
    1 & i = 1, \ldots, m
      \quad \text{(treatment)} \\
    0 & i = m+1, \ldots, N
      \quad \text{(placebo)}
  \end{cases}
\]

\begin{bullets}
\textbf{Three fixed-effect slope parameters describe the run-in, placebo post-randomization, and treatment trajectories under equally spaced observations.}
\begin{itemize}\setlength{\itemsep}{2pt}\setlength{\parskip}{0pt}
\item $\delta$ is the run-in rate of change common to all participants.
\item $\beta$ is the placebo post-randomization rate of change, and $\gamma$ is the treatment effect on that rate.
\item Observations are equally spaced with interval $t$, so $t_j = jt$.
\end{itemize}
\end{bullets}

\begin{rgt}
To be completed by rgt.
\end{rgt}

\begin{orig}
The fixed effects are: \(\delta\), the rate of change during the
run-in period (common to all participants); \(\beta\), the rate
of change during the post-randomization period for the placebo
group; and \(\gamma\), the treatment effect on the rate of
change. We assume equally spaced observations with interval
\(t\), so \(t_j = jt\).

\end{orig}

\begin{bullets}
\textbf{During the common close both arms revert to a shared off-treatment slope, modelled by deactivating the treatment indicator in that phase.}
\begin{itemize}\setlength{\itemsep}{2pt}\setlength{\parskip}{0pt}
\item For $j > J_1$ all participants are off-treatment with slope $\beta + b_i$.
\item This is implemented by setting $g_i = 0$ for every participant during the common close.
\item Equivalently, the treatment indicator is replaced with $g_i \cdot I(1 \le j \le J_1)$ in the model.
\end{itemize}
\end{bullets}

\begin{rgt}
To be completed by rgt.
\end{rgt}

\begin{orig}
For the common close phase (\(j > J_1\)), both treatment and
placebo groups revert to a common off-treatment slope. We
model this by setting \(g_i = 0\) for all participants during
the common close, so that the slope for both groups is \(\beta
+ b_i\) in this phase. This is equivalent to replacing the
treatment indicator with \(g_i \cdot I(1 \le j \le J_1)\) in
model \eqref{eq:model}.

\end{orig}

\subsection{Change-score parameterization}\label{change-score-parameterization}

Following\textsuperscript{\citeproc{ref-frost2008}{2}}, we work with
change scores relative to the randomization visit:
\begin{equation}
  c_{ij} = Y_{ij} - Y_{i0}
  \label{eq:change}
\end{equation}

\begin{bullets}
\textbf{Differencing relative to baseline removes the intercept and aligns the analysis with constrained longitudinal data analysis under compound symmetry.}
\begin{itemize}\setlength{\itemsep}{2pt}\setlength{\parskip}{0pt}
\item The change score eliminates $\alpha$ and the random intercept $a_i$, leaving only two variance components.
\item It is algebraically equivalent to cLDA, so the derived formulas apply to both analyses of run-in designs.
\item By dropping $\sigma_a^2$ it avoids the anticonservative sample sizes that arise from ignoring intercept-slope correlation.
\end{itemize}
\end{bullets}

\begin{rgt}
To be completed by rgt.
\end{rgt}

\begin{orig}
This eliminates the intercept \(\alpha\) and the random
intercept \(a_i\). We note that the change-score approach
is algebraically equivalent to constrained longitudinal
data analysis (cLDA) under compound symmetry, where the
baseline is included in the response vector with a
common-mean constraint across treatment groups\textsuperscript{\citeproc{ref-laird1982}{12},\citeproc{ref-wang2022}{13}}.
The variance formulas derived below therefore apply
to both change-score and cLDA analyses of run-in
designs. Because the change-score eliminates \(a_i\),
the resulting formulas depend on only two variance
components (\(\sigma^2, \sigma_b^2\)) rather than the
three (\(\sigma^2, \sigma_a^2, \sigma_b^2\)) required
by the raw-outcome model.\textsuperscript{\citeproc{ref-zhao2022}{14}} show that
ignoring the intercept-slope correlation in the
three-component model can produce anticonservative
sample sizes; the change-score approach avoids this
issue by construction, at the cost of discarding
information about \(\sigma_a^2\).

\end{orig}

The change-score formulation yields
\begin{equation}
  c_{ij} = \bigl(\delta(1-h_j) + \beta h_j
    + \gamma h_j g_i + b_i\bigr) t_j
    + (w_{ij} - w_{i0})
  \label{eq:change_model}
\end{equation}
for \(j \ne 0\). The change-score vector for participant \(i\)
has dimension \(J = J_0 + J_1 + J_2\), encompassing all three phases.

\subsection{Covariance structure}\label{covariance-structure}

The marginal covariance of the change-score vector is
\begin{equation}
  \Sigma = R + Z G Z'
  \label{eq:sigma}
\end{equation}
where \(G = \sigma_b^2\) (scalar, since we have a single random
slope), \(Z = t \cdot \mathbf{1}_J\) is the \(J \times 1\) random
effects design vector, and \(R\) is the \(J \times J\) residual
covariance of the differenced within-visit residuals
\((w_{ij} - w_{i0})\).

The residual covariance has the structure
\begin{equation}
  R_{jl} = \mathop{\mathrm{Cov}}(w_{ij} - w_{i0},\; w_{il} - w_{i0})
  = \begin{cases}
    2\sigma^2 & \text{if } j = l \\
    \sigma^2  & \text{if } j \ne l
  \end{cases}
  \label{eq:R}
\end{equation}

\begin{bullets}
\textbf{The shared baseline residual induces compound-symmetry correlation among change scores, so the residual covariance depends only on the summed within-visit variance.}
\begin{itemize}\setlength{\itemsep}{2pt}\setlength{\parskip}{0pt}
\item The within-visit residuals are independent across visits with common variance $\sigma^2$.
\item Differencing against the shared baseline $w_{i0}$ produces positive correlation, giving $R = \sigma^2(\mathbf{I}_J + \mathbf{1}_J \mathbf{1}_J')$.
\item Both within-visit components enter $R$ identically, so only their sum is identified from single-scan designs.
\end{itemize}
\end{bullets}

\begin{rgt}
To be completed by rgt.
\end{rgt}

\begin{orig}
since \(w_{ij} = u_{ij} + \varepsilon_{ij}\) are independent
across visits with common variance
\(\sigma^2 = \sigma_u^2 + \sigma_\varepsilon^2\), and the
shared baseline residual \(w_{i0}\) induces positive
correlation among all change scores. Thus
\(R = \sigma^2(\mathbf{I}_J + \mathbf{1}_J \mathbf{1}_J')\).
Both within-visit components enter \(R\) identically because
\(u_{i0}\) and \(\varepsilon_{i0}\) are differenced out of every
change score in the same way; this is why only their sum
\(\sigma^2\) is identified from single-scan designs.

\end{orig}

\section{Variance of the Treatment Effect Estimator}\label{variance-of-the-treatment-effect-estimator}

\subsection{GLS estimation and the Woodbury identity}\label{gls-estimation-and-the-woodbury-identity}

The generalized least squares estimator of the fixed effects
is
\begin{equation}
  \hat{\boldsymbol{\beta}} = (X' \Sigma^{-1} X)^{-1} X' \Sigma^{-1}
  \mathbf{c}
  \label{eq:gls}
\end{equation}

\begin{bullets}
\textbf{The estimand of interest is the treatment-effect variance, a single diagonal element of the inverted GLS information matrix.}
\begin{itemize}\setlength{\itemsep}{2pt}\setlength{\parskip}{0pt}
\item The GLS estimator has covariance $(X' \Sigma^{-1} X)^{-1}$.
\item $X$ is the fixed-effects design matrix and $\mathbf{c}$ is the stacked change-score vector.
\item The target $\mathop{\mathrm{Var}}(\hat{\gamma})$ is the diagonal entry corresponding to $\gamma$.
\end{itemize}
\end{bullets}

\begin{rgt}
To be completed by rgt.
\end{rgt}

\begin{orig}
with \(\mathop{\mathrm{Var}}(\hat{\boldsymbol{\beta}}) = (X' \Sigma^{-1} X)^{-1}\),
where \(X\) is the fixed effects design matrix and \(\mathbf{c}\)
is the stacked change-score vector. Our target is
\(\mathop{\mathrm{Var}}(\hat{\gamma})\), the variance of the treatment effect
estimator, which is the diagonal element of
\((X' \Sigma^{-1} X)^{-1}\) corresponding to \(\gamma\).

\end{orig}

The key computational tool is the Woodbury matrix identity.
Since \(\Sigma = R + Z G Z'\) with scalar \(G = \sigma_b^2\):
\begin{equation}
  \Sigma^{-1} = R^{-1} - R^{-1} Z
    \bigl(G^{-1} + Z' R^{-1} Z\bigr)^{-1}
    Z' R^{-1}
  \label{eq:woodbury}
\end{equation}

This avoids direct inversion of the \(J \times J\) matrix
\(\Sigma\) and enables closed-form results because \(R^{-1}\)
has a known simple structure.

From \(R = \sigma^2(\mathbf{I}_J + \mathbf{1}_J
\mathbf{1}_J')\), the Sherman-Morrison formula gives
\begin{equation}
  R^{-1} = \frac{1}{\sigma^2}\left(
    \mathbf{I}_J - \frac{1}{J+1}
    \mathbf{1}_J \mathbf{1}_J'
  \right)
  \label{eq:Rinv}
\end{equation}

With \(Z = t \cdot \mathbf{1}_J\), we obtain
\begin{equation}
  Z' R^{-1} Z = \frac{t^2}{\sigma^2}
    \left(J - \frac{J^2}{J+1}\right)
  = \frac{J\, t^2}{\sigma^2(J+1)}
  \label{eq:ZRZ}
\end{equation}
and
\begin{equation}
  \bigl(G^{-1} + Z' R^{-1} Z\bigr)^{-1}
  = \frac{\sigma^2 \sigma_b^2 (J+1)}
    {\sigma^2(J+1) + J\, t^2 \sigma_b^2}
  \label{eq:H}
\end{equation}

Let \(a = \sigma^2\) and \(b = \sigma_b^2\) for conciseness.
Define
\begin{equation}
  \eta_J = \frac{a(J+1) + J\, t^2 b}{a(J+1)}
  = 1 + \frac{J\, t^2 b}{a(J+1)}
  \label{eq:eta}
\end{equation}
so that \((G^{-1} + Z'R^{-1}Z)^{-1} = b / \eta_J\).

\subsubsection{Per-participant information, compact form}\label{per-participant-information-compact-form}

Let \(X_i\), \(Z_i\), and \(R_i\) denote the per-participant fixed
effects design, random effects design, and residual covariance,
and define the scalar
\begin{equation}
  \mathcal{J}_i
  = Z_i' R_i^{-1} Z_i + G^{-1}
  \label{eq:Ji}
\end{equation}
By (\ref{eq:woodbury}), the per-participant contribution to
the GLS information matrix is
\begin{equation}
  X_i' \Sigma_i^{-1} X_i
  = X_i' R_i^{-1} X_i
    - \frac{(X_i' R_i^{-1} Z_i)(Z_i' R_i^{-1} X_i)}
           {\mathcal{J}_i}
  \label{eq:perpt_info}
\end{equation}

\begin{bullets}
\textbf{The additive per-participant information form has two properties exploited later: a scalar rank-one correction and accommodation of heterogeneous designs.}
\begin{itemize}\setlength{\itemsep}{2pt}\setlength{\parskip}{0pt}
\item The total information is the sum of per-participant contributions, with covariance given by its inverse.
\item Under a single random slope, $\mathcal{J}_i$ is a scalar, making the correction a rank-one adjustment of the OLS-style information.
\item Because contributions add, varying $J_0$, $J_2$, or dropout patterns are handled by summing over strata.
\end{itemize}
\end{bullets}

\begin{rgt}
To be completed by rgt.
\end{rgt}

\begin{orig}
and the total information is \(X' \Sigma^{-1} X = \sum_i
X_i' \Sigma_i^{-1} X_i\), with \(\mathop{\mathrm{Var}}(\hat{\boldsymbol{\beta}})
= (X' \Sigma^{-1} X)^{-1}\). Two features of this form are
used later: (i) under a single random slope, \(\mathcal{J}_i\)
is a scalar so (\ref{eq:perpt_info}) is a rank-one correction
of the OLS-style information \(X_i' R_i^{-1} X_i\); (ii) because
per-participant contributions add, heterogeneous designs
(varying \(J_0\), \(J_2\), or dropout patterns) are accommodated
by summing (\ref{eq:perpt_info}) over the relevant strata.

\end{orig}

\subsection{\texorpdfstring{Case 1: Standard design, no run-in (\(J_0=0\), \(J_1=2\), \(J_2=0\))}{Case 1: Standard design, no run-in (J\_0=0, J\_1=2, J\_2=0)}}\label{case-1-standard-design-no-run-in-j_00-j_12-j_20}

\begin{bullets}
\textbf{The no-run-in case reproduces established results and serves as a verification baseline for the framework.}
\begin{itemize}\setlength{\itemsep}{2pt}\setlength{\parskip}{0pt}
\item It reproduces Frost and colleagues, Section 3.1.1.
\item It is the $J_0 = 0$ baseline of the random-coefficient power formulae of Lu and Hu and colleagues.
\item It thereby provides a verification checkpoint for the present derivation.
\end{itemize}
\end{bullets}

\begin{rgt}
To be completed by rgt.
\end{rgt}

\begin{orig}
This case reproduces the result in\textsuperscript{\citeproc{ref-frost2008}{2}}, Section 3.1.1, and
is the \(J_0 = 0\) baseline of the random coefficient regression
power formulae of\textsuperscript{\citeproc{ref-lu2009}{4}} and\textsuperscript{\citeproc{ref-hu2021}{6}}. It serves as
verification for our framework.

\end{orig}

\begin{bullets}
\textbf{Without a run-in phase the design reduces to two post-randomization change scores and a two-parameter fixed-effects model.}
\begin{itemize}\setlength{\itemsep}{2pt}\setlength{\parskip}{0pt}
\item There are $J = 2$ change scores at post-randomization times $t$ and $2t$.
\item The absence of run-in removes the $\delta$ parameter from the model.
\item The remaining fixed effects are the placebo slope $\beta$ and the treatment effect $\gamma$.
\end{itemize}
\end{bullets}

\begin{rgt}
To be completed by rgt.
\end{rgt}

\begin{orig}
With \(J_0=0\) and \(J_1=2\), there are \(J = 2\) change scores
\((c_{i,1},\, c_{i,2})\) at post-randomization times \((t,\,
2t)\). Since there is no run-in, the model has two fixed
effects: the placebo slope \(\beta\) and the treatment effect
\(\gamma\).

\end{orig}

The design matrices are:
\[
  X_i^{(\text{trt})} = t\begin{pmatrix}
    1 & 1 \\ 2 & 2
  \end{pmatrix}, \qquad
  X_i^{(\text{plc})} = t\begin{pmatrix}
    1 & 0 \\ 2 & 0
  \end{pmatrix}
\]

The residual covariance is
\(R = a \begin{pmatrix} 2 & 1 \\ 1 & 2 \end{pmatrix}\)
with \(R^{-1} = \frac{1}{3a}
\begin{pmatrix} 2 & -1 \\ -1 & 2 \end{pmatrix}\),
and the random effects design is \(Z = t(1,\, 2)'\).

Applying the Woodbury identity with \(J = 2\):
\begin{equation}
  Z'R^{-1}Z = \frac{t^2}{3a}(1,\,2)
    \begin{pmatrix}2&-1\\-1&2\end{pmatrix}
    \begin{pmatrix}1\\2\end{pmatrix}
  = \frac{t^2}{3a}(1,\,2)(0,\,3)' = \frac{6t^2}{3a}
  = \frac{2t^2}{a}
\end{equation}
\begin{equation}
  (G^{-1} + Z'R^{-1}Z)^{-1}
  = \frac{ab}{a + 2bt^2}
\end{equation}

After summing over one treatment and one placebo participant,
the information matrix is
\begin{equation}
  X'\Sigma^{-1}X
  = \frac{2t^2}{a + 2bt^2}
    \begin{pmatrix} 2 & 1 \\ 1 & 1 \end{pmatrix}
\end{equation}
and its inverse is
\begin{equation}
  \mathop{\mathrm{Var}}(\hat{\boldsymbol{\beta}})
  = \frac{a + 2bt^2}{2t^2}
    \begin{pmatrix} 1 & -1 \\ -1 & 2 \end{pmatrix}
  \label{eq:var_frost}
\end{equation}

The variance of the treatment effect estimator (per
participant pair) is the \((2,2)\) element:
\begin{equation}
  \boxed{
    \mathop{\mathrm{Var}}(\hat{\gamma})\big|_{J_0=0,J_1=2}
    = \frac{a + 2bt^2}{t^2}
    = \frac{\sigma^2 + 2\sigma_b^2 t^2}{t^2}
  }
  \label{eq:var_r0k2}
\end{equation}

\begin{bullets}
\textbf{The recovered variance matches Frost and supplies the per-group baseline for all run-in comparisons.}
\begin{itemize}\setlength{\itemsep}{2pt}\setlength{\parskip}{0pt}
\item The expression agrees with Frost and colleagues, Section 3.1.1.
\item For $m$ participants per group the variance scales as $(\sigma^2 + 2\sigma_b^2 t^2)/(m t^2)$.
\item This quantity is the reference against which run-in designs are evaluated.
\end{itemize}
\end{bullets}

\begin{rgt}
To be completed by rgt.
\end{rgt}

\begin{orig}
This agrees with Frost et al.~(2008), Section 3.1.1. For
\(m\) participants per group:
\(\mathop{\mathrm{Var}}(\hat{\gamma}) = (\sigma^2 + 2\sigma_b^2 t^2) /
(m t^2)\). This serves as the baseline against which the
run-in designs are compared.

\end{orig}

\subsection{\texorpdfstring{Case 2: Two run-in observations (\(J_0=2\), \(J_1=1\), \(J_2=0\))}{Case 2: Two run-in observations (J\_0=2, J\_1=1, J\_2=0)}}\label{case-2-two-run-in-observations-j_02-j_11-j_20}

\begin{bullets}
\textbf{Two run-in observations yield three change scores and a design matrix that distinguishes the run-in slope from the post-randomization slopes.}
\begin{itemize}\setlength{\itemsep}{2pt}\setlength{\parskip}{0pt}
\item There are $J = 3$ change scores spanning the pre- and post-randomization times.
\item The fixed-effects model now includes the run-in slope $\delta$ alongside $\beta$ and $\gamma$.
\item The design matrix separates the run-in contribution from the post-randomization contribution.
\end{itemize}
\end{bullets}

\begin{rgt}
To be completed by rgt.
\end{rgt}

\begin{orig}
We now have \(J = 3\) change scores \((c_{i,-2},\, c_{i,-1},\,
c_{i,1})\) at times \((-2t,\, -t,\, 0,\, t)\). The fixed effects
design matrix separates the run-in slope \(\delta\) from the
post-randomization slopes:

\end{orig}

\[
  X_i^{(\text{trt})} = \begin{pmatrix}
    -2t & 0 & 0 \\ -t & 0 & 0 \\ 0 & t & t
  \end{pmatrix}, \qquad
  X_i^{(\text{plc})} = \begin{pmatrix}
    -2t & 0 & 0 \\ -t & 0 & 0 \\ 0 & t & 0
  \end{pmatrix}
\]

\begin{bullets}
\textbf{Because the random slope acts uniformly across visits, the random-effects design vector collapses to the scaled time distances from baseline.}
\begin{itemize}\setlength{\itemsep}{2pt}\setlength{\parskip}{0pt}
\item The random slope $b_i$ applies to all time points alike.
\item The random-effects design therefore reduces to $Z = t(-2,\, -1,\, 1)'$.
\item The run-in change scores are differences of the pre-randomization outcomes from baseline.
\end{itemize}
\end{bullets}

\begin{rgt}
To be completed by rgt.
\end{rgt}

\begin{orig}
The random effects design is \(Z = t(1,\, 1,\, 1)' \cdot
(-2,\, -1,\, 1)' \cdot b_i\). However, since the random
slope applies uniformly to all time points, \(Z = t(-2,\, -1,
\, 1)'\) for the random effect \(b_i\). We note that the
change scores at times \(-2t\) and \(-t\) relative to baseline
at \(0\) yield \(c_{i,-2} = Y_{i,-2} - Y_{i,0}\) and
\(c_{i,-1} = Y_{i,-1} - Y_{i,0}\).

\end{orig}

More carefully, under model \eqref{eq:model}:
\begin{align}
  c_{i,-2} &= -2(\delta + b_i)\,t + (w_{i,-2} - w_{i,0})
    \notag \\
  c_{i,-1} &= -(\delta + b_i)\,t + (w_{i,-1} - w_{i,0})
    \notag \\
  c_{i,1}  &= (\beta + \gamma g_i + b_i)\,t
    + (w_{i,1} - w_{i,0})
  \label{eq:c_r2}
\end{align}

Thus the design matrices for the fixed effects
\((\delta,\, \beta,\, \gamma)'\) are:
\[
  X_i^{(\text{trt})} = t\begin{pmatrix}
    -2 & 0 & 0 \\ -1 & 0 & 0 \\ 0 & 1 & 1
  \end{pmatrix}, \qquad
  X_i^{(\text{plc})} = t\begin{pmatrix}
    -2 & 0 & 0 \\ -1 & 0 & 0 \\ 0 & 1 & 0
  \end{pmatrix}
\]
and the random effects design is
\(Z_i = t(-2,\, -1,\, 1)'\).

The residual covariance is
\begin{equation}
  R = a\begin{pmatrix}
    2 & 1 & 1 \\ 1 & 2 & 1 \\ 1 & 1 & 2
  \end{pmatrix}
  = a(\mathbf{I}_3 + \mathbf{1}_3\mathbf{1}_3')
\end{equation}
with inverse
\begin{equation}
  R^{-1} = \frac{1}{a}\left(\mathbf{I}_3
    - \frac{1}{4}\mathbf{1}_3\mathbf{1}_3'\right)
  = \frac{1}{4a}\begin{pmatrix}
    3 & -1 & -1 \\ -1 & 3 & -1 \\ -1 & -1 & 3
  \end{pmatrix}
\end{equation}

Now we compute the Woodbury components. With
\(Z = t(-2,\, -1,\, 1)'\):
\begin{equation}
  Z'R^{-1}Z = \frac{t^2}{4a}\bigl(
    3 \cdot 4 + 3 \cdot 1 + 3 \cdot 1
    - 2(-1)(-2)(-1) - 2(-1)(-2)(1)
    - 2(-1)(-1)(1)\bigr)
\end{equation}

Computing directly:
\(Z'R^{-1}Z = \frac{t^2}{4a}(-2,\,-1,\,1)
\begin{pmatrix} 3&-1&-1 \\ -1&3&-1 \\ -1&-1&3
\end{pmatrix}
(-2,\,-1,\,1)'\)

\begin{align}
  R^{-1}Z &= \frac{t}{4a}\begin{pmatrix}
    3(-2)+(-1)(-1)+(-1)(1) \\
    (-1)(-2)+3(-1)+(-1)(1) \\
    (-1)(-2)+(-1)(-1)+3(1)
  \end{pmatrix}
  = \frac{t}{4a}\begin{pmatrix}
    -6 \\ -2 \\ 6
  \end{pmatrix}
  \notag
\end{align}
so \(Z'R^{-1}Z = \frac{t^2}{4a}
[(-2)(-6)+(-1)(-2)+(1)(6)] = \frac{t^2}{4a}(12+2+6)
= \frac{20t^2}{4a} = \frac{5t^2}{a}\).

Alternatively, we note \(Z = t \cdot \mathbf{d}\) where
\(\mathbf{d} = (-2,-1,1)'\), so \(Z'R^{-1}Z =
\frac{t^2}{a}(\mathbf{d}'\mathbf{d} -
\frac{1}{4}(\mathbf{1}'\mathbf{d})^2) =
\frac{t^2}{a}(6 - \frac{4}{4}) = \frac{5t^2}{a}\).

Thus
\begin{equation}
  H \equiv (G^{-1}+Z'R^{-1}Z)^{-1}
  = \frac{ab}{a + 5bt^2}
\end{equation}

For the information matrix, we compute
\(X'R^{-1}X\) summed over treatment and placebo groups.
For a single treatment participant:
\[
  X_i^{(\text{trt})\prime} R^{-1} X_i^{(\text{trt})}
  = \frac{t^2}{4a}\begin{pmatrix}
    -2 & -1 & 0 \\ 0 & 0 & 1 \\ 0 & 0 & 1
  \end{pmatrix}
  \begin{pmatrix}
    3&-1&-1\\-1&3&-1\\-1&-1&3
  \end{pmatrix}
  \begin{pmatrix}
    -2&0&0\\-1&0&0\\0&1&1
  \end{pmatrix}
\]

Computing the product:
\begin{equation}
  X_i^{(\text{trt})\prime} R^{-1} X_i^{(\text{trt})}
  = \frac{t^2}{4a}\begin{pmatrix}
    11 & 3 & 3 \\ 3 & 3 & 3 \\ 3 & 3 & 3
  \end{pmatrix}
\end{equation}

For placebo (whose \(\gamma\) column is zero):
\begin{equation}
  X_i^{(\text{plc})\prime} R^{-1} X_i^{(\text{plc})}
  = \frac{t^2}{4a}\begin{pmatrix}
    11 & 3 & 0 \\ 3 & 3 & 0 \\ 0 & 0 & 0
  \end{pmatrix}
\end{equation}

Summing (one participant per group for the derivation, then
multiplying by \(m\)):
\begin{equation}
  X'R^{-1}X = \frac{t^2}{4a}\begin{pmatrix}
    22 & 6 & 3 \\ 6 & 6 & 3 \\ 3 & 3 & 3
  \end{pmatrix}
\end{equation}

For the Woodbury correction, we need
\(X'R^{-1}Z\). For treatment:
\[
  X_i^{(\text{trt})\prime}R^{-1}Z
  = \frac{t^2}{4a}\begin{pmatrix}-2&-1&0\\0&0&1\\0&0&1
  \end{pmatrix}
  \begin{pmatrix}-6\\-2\\6\end{pmatrix}
  = \frac{t^2}{4a}\begin{pmatrix}14\\6\\6\end{pmatrix}
\]

For placebo:
\[
  X_i^{(\text{plc})\prime}R^{-1}Z
  = \frac{t^2}{4a}\begin{pmatrix}-2&-1&0\\0&0&1\\0&0&0
  \end{pmatrix}
  \begin{pmatrix}-6\\-2\\6\end{pmatrix}
  = \frac{t^2}{4a}\begin{pmatrix}14\\6\\0\end{pmatrix}
\]

The Woodbury correction requires the per-subject outer
products (not the outer product of the sum):
\begin{align}
  W &= \sum_{g} X_i^{(g)\prime} R^{-1} Z
    \cdot H \cdot
    Z' R^{-1} X_i^{(g)} \notag \\
  &= \frac{bt^4}{16a(a+5bt^2)}
    \left(
      \begin{pmatrix}14\\6\\6\end{pmatrix}
      (14,6,6)
      + \begin{pmatrix}14\\6\\0\end{pmatrix}
      (14,6,0)
    \right) \notag \\
  &= \frac{bt^4}{16a(a+5bt^2)}
    \begin{pmatrix}
      392 & 168 & 84 \\
      168 & 72 & 36 \\
      84 & 36 & 36
    \end{pmatrix}
\end{align}

Therefore
\begin{equation}
  X'\Sigma^{-1}X = X'R^{-1}X - W
\end{equation}

\begin{bullets}
\textbf{The treatment-effect variance is the relevant diagonal element of the inverted information matrix, obtained numerically and checked against the closed form.}
\begin{itemize}\setlength{\itemsep}{2pt}\setlength{\parskip}{0pt}
\item $\mathop{\mathrm{Var}}(\hat{\gamma})$ is the $(3,3)$ element of $(X'\Sigma^{-1}X)^{-1}$.
\item It is evaluated numerically in the accompanying R code.
\item The numerical value is verified against the derived closed-form result.
\end{itemize}
\end{bullets}

\begin{rgt}
To be completed by rgt.
\end{rgt}

\begin{orig}
To find \(\mathop{\mathrm{Var}}(\hat{\gamma})\), we need the \((3,3)\) element
of \((X'\Sigma^{-1}X)^{-1}\). We compute this numerically in
the R code below and verify against the closed-form result.

\end{orig}

\begin{bullets}
\textbf{Completing the cofactor and determinant algebra yields the closed-form treatment-effect variance for the two-run-in design.}
\begin{itemize}\setlength{\itemsep}{2pt}\setlength{\parskip}{0pt}
\item Define the auxiliary quantity $D = \sigma^2 + 5\sigma_b^2 t^2$.
\item The $(3,3)$ cofactor simplifies to $6t^4/(aD)$.
\item Dividing the cofactor by the determinant gives the boxed variance below.
\end{itemize}
\end{bullets}

\begin{rgt}
To be completed by rgt.
\end{rgt}

\begin{orig}
Carrying through the algebra, define \(D = a + 5bt^2 =
\sigma^2 + 5\sigma_b^2 t^2\) and write
\(X'\Sigma^{-1}X = \frac{t^2}{16aD}\,A\). The matrix \(A\) has
\((3,3)\) minor \(1536\,aD\) and determinant
\(9216\,aD\,(a+2bt^2)\). Hence the \((3,3)\) cofactor of
\(X'\Sigma^{-1}X\) is
\(1536aD \cdot t^4/(16aD)^2 = 6t^4/(aD)\), and its
determinant is
\(9216\,aD\,(a+2bt^2)\cdot t^6/(16aD)^3\). The ratio
gives:

\end{orig}

\begin{equation}
  \boxed{
    \mathop{\mathrm{Var}}(\hat{\gamma})\big|_{J_0=2,J_1=1}
    = \frac{8\sigma^2(\sigma^2 + 5\sigma_b^2 t^2)}
           {3t^2(\sigma^2 + 2\sigma_b^2 t^2)}
  }
  \label{eq:var_r2k1}
\end{equation}

\begin{bullets}
\textbf{The run-in variance differs structurally from the Frost baseline by involving the residual variance in both numerator and denominator.}
\begin{itemize}\setlength{\itemsep}{2pt}\setlength{\parskip}{0pt}
\item In the Frost formula $\sigma^2$ appears only additively in the numerator.
\item Here $\sigma^2$ enters both the numerator and the denominator.
\item This reflects the extra information extracted from pre-randomization observations.
\end{itemize}
\end{bullets}

\begin{rgt}
To be completed by rgt.
\end{rgt}

\begin{orig}
Unlike the Frost formula \eqref{eq:var_r0k2}, the run-in
variance depends on \(\sigma^2\) both in the numerator and
denominator, reflecting the additional information extracted
from pre-randomization observations.

\end{orig}

\subsection{\texorpdfstring{General formula: \(J_0\) run-in observations (\(J_2=0\))}{General formula: J\_0 run-in observations (J\_2=0)}}\label{general-formula-j_0-run-in-observations-j_20}

\begin{bullets}
\textbf{The derivation now generalizes to an arbitrary number of run-in observations with no common close.}
\begin{itemize}\setlength{\itemsep}{2pt}\setlength{\parskip}{0pt}
\item The design has arbitrary $J_0$ run-in and $J_1$ post-randomization observations.
\item The common close is absent, so $J_2 = 0$.
\item The total number of change scores is $J = J_0 + J_1$.
\end{itemize}
\end{bullets}

\begin{rgt}
To be completed by rgt.
\end{rgt}

\begin{orig}
We now derive the variance for arbitrary \(J_0\) run-in
observations with \(J_1\) post-randomization observations and
no common close (\(J_2 = 0\)). The total number of change scores
is \(J = J_0 + J_1\).

\end{orig}

The change scores are
\begin{align}
  c_{i,-s} &= -s(\delta + b_i)\,t + (w_{i,-s} - w_{i0})
    & s &= 1,\ldots,J_0 \notag \\
  c_{i,l}  &= l(\beta + \gamma g_i + b_i)\,t
    + (w_{i,l} - w_{i0})
    & l &= 1,\ldots,J_1
\end{align}

\begin{bullets}
\textbf{The design matrix is built from run-in and post-randomization time vectors mapping to the three fixed-effect columns.}
\begin{itemize}\setlength{\itemsep}{2pt}\setlength{\parskip}{0pt}
\item The columns correspond to the parameters $\delta$, $\beta$, and $\gamma$.
\item $\mathbf{d}_{J_0}$ collects the negative run-in times scaled by $t$.
\item $\mathbf{d}_{J_1}$ collects the positive post-randomization times scaled by $t$.
\end{itemize}
\end{bullets}

\begin{rgt}
To be completed by rgt.
\end{rgt}

\begin{orig}
The fixed effects design matrix has columns for
\((\delta, \beta, \gamma)\). Let \(\mathbf{d}_{J_0} = (-J_0, -J_0+1,
\ldots, -1)' \cdot t\) be the run-in time vector and
\(\mathbf{d}_{J_1} = (1, 2, \ldots, J_1)' \cdot t\) be the
post-randomization time vector. Then:

\end{orig}

\[
  X_i^{(\text{trt})} = \begin{pmatrix}
    \mathbf{d}_{J_0} & \mathbf{0} & \mathbf{0} \\
    \mathbf{0}   & \mathbf{d}_{J_1} & \mathbf{d}_{J_1}
  \end{pmatrix}, \qquad
  X_i^{(\text{plc})} = \begin{pmatrix}
    \mathbf{d}_{J_0} & \mathbf{0} & \mathbf{0} \\
    \mathbf{0}   & \mathbf{d}_{J_1} & \mathbf{0}
  \end{pmatrix}
\]

The random effects design vector is
\(Z_i = (\mathbf{d}_{J_0}',\, \mathbf{d}_{J_1}')'\), i.e., the full
vector of time distances from baseline.

The residual covariance remains \(R = a(\mathbf{I}_J +
\mathbf{1}_J\mathbf{1}_J')\) with inverse given by
\eqref{eq:Rinv}.

\subsubsection{\texorpdfstring{Computing \(X'\Sigma^{-1}X\)}{Computing X\textquotesingle\textbackslash Sigma\^{}\{-1\}X}}\label{computing-xsigma-1x}

\begin{bullets}
\textbf{The Woodbury computation is organized around the sum-of-squares of the run-in and post-randomization time indices.}
\begin{itemize}\setlength{\itemsep}{2pt}\setlength{\parskip}{0pt}
\item The derivation proceeds via the Woodbury identity.
\item $S_{J_0}$ is the sum of squared run-in indices in closed form.
\item $S_{J_1}$ is the corresponding sum over post-randomization indices.
\end{itemize}
\end{bullets}

\begin{rgt}
To be completed by rgt.
\end{rgt}

\begin{orig}
We proceed through the Woodbury identity. Define
\(S_{J_0} = \sum_{s=1}^{J_0} s^2 = J_0(J_0+1)(2J_0+1)/6\) and
\(S_{J_1} = \sum_{l=1}^{J_1} l^2 = J_1(J_1+1)(2J_1+1)/6\).

\end{orig}

The key intermediate quantities are:

\textbf{\(Z'R^{-1}Z\):} With \(Z = t \cdot \mathbf{d}\) where
\(\mathbf{d} = (-J_0,\ldots,-1,1,\ldots,J_1)'\):
\begin{equation}
  Z'R^{-1}Z = \frac{t^2}{a}\left[
    (S_{J_0} + S_{J_1})
    - \frac{(\sum d_j)^2}{J+1}
  \right]
  \label{eq:ZRZ_general}
\end{equation}
where \(\sum d_j = -J_0(J_0+1)/2 + J_1(J_1+1)/2\).

\textbf{\((G^{-1} + Z'R^{-1}Z)^{-1}\):} From \eqref{eq:ZRZ_general}:
\[
  H = \frac{ab}{a + t^2 b \cdot Z'R^{-1}Z \cdot a/t^2}
\]

\textbf{\(X'R^{-1}X\):} This \(3 \times 3\) matrix has elements
computed from the block structure of \(X\). Letting
\(D_{J_0} = \sum_{s=1}^{J_0} s = J_0(J_0+1)/2\) and
\(D_{J_1} = \sum_{l=1}^{J_1} l = J_1(J_1+1)/2\), and using
\(R^{-1} = \frac{1}{a}(\mathbf{I}_J -
\frac{1}{J+1}\mathbf{1}\mathbf{1}')\), the information matrix
summed over one treatment and one placebo participant is:

\begin{equation}
  X'R^{-1}X = \frac{t^2}{a}\begin{pmatrix}
    2\!\left(S_{J_0} - \frac{D_{J_0}^2}{J+1}\right) &
    \frac{2D_{J_0} D_{J_1}}{J+1} &
    \frac{D_{J_0} D_{J_1}}{J+1} \\[4pt]
    \frac{2D_{J_0} D_{J_1}}{J+1} &
    2\!\left(S_{J_1} - \frac{D_{J_1}^2}{J+1}\right) &
    S_{J_1} - \frac{D_{J_1}^2}{J+1} \\[4pt]
    \frac{D_{J_0} D_{J_1}}{J+1} &
    S_{J_1} - \frac{D_{J_1}^2}{J+1} &
    S_{J_1} - \frac{D_{J_1}^2}{J+1}
  \end{pmatrix}
  \label{eq:XRX_general}
\end{equation}

\begin{bullets}
\textbf{The structure of the information matrix reflects which parameters both groups inform and the opposing signs of the two phase time-sums.}
\begin{itemize}\setlength{\itemsep}{2pt}\setlength{\parskip}{0pt}
\item Entries informed by both groups carry a factor of 2.
\item The $\gamma$-related entries lack this factor because the placebo $\gamma$ column is zero.
\item The positive $(\delta, \beta)$ cross-term arises from the opposite signs of the run-in and post-randomization time sums.
\end{itemize}
\end{bullets}

\begin{rgt}
To be completed by rgt.
\end{rgt}

\begin{orig}
The factor of 2 appears on entries involving both groups,
while \(\gamma\)-related entries lack this factor since
the \(\gamma\) column is zero for the placebo group. The
positive sign on the \((\delta, \beta)\) cross-term arises
because the run-in and post-randomization time sums have
opposite signs (\(\sum d_{J_0} < 0\), \(\sum d_{J_1} > 0\)).

\end{orig}

\textbf{Woodbury correction:} The correction is the sum of
per-subject rank-one outer products:
\[
  W = \sum_{g \in \{\text{trt},\text{plc}\}}
    X_i^{(g)\prime} R^{-1} Z \cdot H \cdot
    Z' R^{-1} X_i^{(g)}
\]
which is generally rank 2 (not rank 1), so
\(X'\Sigma^{-1}X = X'R^{-1}X - W\).

\subsubsection{Numerical implementation}\label{numerical-implementation}

\begin{bullets}
\textbf{For general run-in and post-randomization counts the variance has no compact closed form and must be evaluated as a complex ratio.}
\begin{itemize}\setlength{\itemsep}{2pt}\setlength{\parskip}{0pt}
\item No simple closed-form expression exists for arbitrary $J_0$ and $J_1$.
\item The variance depends on the sum-of-squares and sum quantities together with the dimension $J$.
\item The Woodbury correction interacts with the information matrix to produce algebraically complex expressions.
\end{itemize}
\end{bullets}

\begin{rgt}
To be completed by rgt.
\end{rgt}

\begin{orig}
The variance \(\mathop{\mathrm{Var}}(\hat{\gamma})\) for general \(J_0\) and \(J_1\)
does not admit a simple closed-form expression. The
formula depends on \(\sigma^2\) and \(\sigma_b^2\) in a
ratio involving \(S_{J_0}\), \(S_{J_1}\), \(D_{J_0}\), \(D_{J_1}\), and \(J\),
and the interaction between the Woodbury correction and the
information matrix produces algebraically complex
expressions.

\end{orig}

\begin{bullets}
\textbf{A matrix-based R implementation evaluates the exact variance for any configuration, with the closed-form cases as checkpoints.}
\begin{itemize}\setlength{\itemsep}{2pt}\setlength{\parskip}{0pt}
\item The computation uses the matrix Woodbury identity in R code.
\item It returns the exact variance for any $(J_0, J_1, J_2)$ configuration.
\item The $J_0=0$ and $J_0=2$ closed forms serve as verification checkpoints.
\end{itemize}
\end{bullets}

\begin{rgt}
To be completed by rgt.
\end{rgt}

\begin{orig}
We therefore implement the computation via the matrix-based
Woodbury identity in R code, which evaluates the exact
variance for any configuration \((J_0, J_1, J_2)\). The
closed-form expressions for \(J_0=0\) (\ref{eq:var_r0k2}) and
\(J_0=2\) (\ref{eq:var_r2k1}) serve as verification checkpoints.

\end{orig}

\subsubsection{Remark: recovery of the Edland and Diggle formulae}\label{remark-recovery-of-the-edland-and-diggle-formulae}

\begin{bullets}
\textbf{Setting the run-in count to zero collapses the information matrix to a two-parameter block and yields the slope-comparison variance.}
\begin{itemize}\setlength{\itemsep}{2pt}\setlength{\parskip}{0pt}
\item With $J_0 = 0$ the $\delta$ column vanishes, leaving a $2 \times 2$ block on $(\beta, \gamma)$.
\item Sherman-Morrison cancellations collapse the Woodbury correction against the remaining structure.
\item The result is the per-arm variance of the treatment-effect estimator for $n$ participants.
\end{itemize}
\end{bullets}

\begin{rgt}
To be completed by rgt.
\end{rgt}

\begin{orig}
At \(J_0 = 0\), the \(\delta\) column of \(X\) is empty and the
information matrix reduces to a \(2 \times 2\) block on
\((\beta, \gamma)\). Carrying the Woodbury algebra through, the
Sherman-Morrison cancellations collapse the Woodbury correction
against the \((2,1; 1,1)\) structure of \(X'R^{-1}X\), and the
variance of the treatment effect estimator for \(n\) participants
per arm becomes

\end{orig}

\begin{equation}
  \mathop{\mathrm{Var}}(\hat{\gamma})\big|_{J_0=0,J_2=0}
  = \frac{2}{n}\!\left[\sigma_b^2 + \frac{\sigma^2}{S_{xx}}\right],
  \qquad
  S_{xx} = \sum_{j=0}^{J_1}(t_j - \bar t)^2,
  \label{eq:var_edland}
\end{equation}

\begin{bullets}
\textbf{The boundary case recovers the Edland, Diggle, and Frost formulae, confirming that the general result embeds the longpower routines as strict special cases.}
\begin{itemize}\setlength{\itemsep}{2pt}\setlength{\parskip}{0pt}
\item The expression is the Ard-Edland slope-comparison formula implemented in the longpower edland routine.
\item Setting the slope variance to zero gives compound symmetry and recovers the Diggle formula.
\item The Frost formula is the two-visit specialization, so the general result quantifies gains accruing only from run-in or common-close observations.
\end{itemize}
\end{bullets}

\begin{rgt}
To be completed by rgt.
\end{rgt}

\begin{orig}
where the baseline visit at \(t_0 = 0\) enters through the
change-score transformation. Equation (\ref{eq:var_edland}) is
the slope-comparison formula of\textsuperscript{\citeproc{ref-ard2011}{8}}
implemented in \texttt{longpower::edland.linear.power()}. Setting
additionally \(\sigma_b^2 = 0\) reduces \(\Sigma\) to compound
symmetry and recovers the formula of\textsuperscript{\citeproc{ref-diggle2002}{9}}
used by \texttt{longpower::diggle.linear.power()}. The Frost formula
(\ref{eq:var_r0k2}) is in turn the \(J_1 = 2\) specialization of
(\ref{eq:var_edland}), since \(S_{xx} = 2t^2\) at that setting.
The general-\((J_0, J_1, J_2)\) result derived in the next
subsection therefore embeds both \texttt{longpower} closed-form
routines as strict boundary cases, with the efficiency gains
quantified in Section 5 accruing entirely from the information
contributed by \(J_0 > 0\) or \(J_2 > 0\).

\end{orig}

\subsection{\texorpdfstring{General formula with common close (\(J_2 \ge 0\))}{General formula with common close (J\_2 \textbackslash ge 0)}}\label{general-formula-with-common-close-j_2-ge-0}

\begin{bullets}
\textbf{The final extension adds common-close observations after the treatment period to the change-score vector.}
\begin{itemize}\setlength{\itemsep}{2pt}\setlength{\parskip}{0pt}
\item $J_2$ common-close observations follow the treatment period.
\item All participants are off-treatment during this phase.
\item The total dimension becomes $J = J_0 + J_1 + J_2$.
\end{itemize}
\end{bullets}

\begin{rgt}
To be completed by rgt.
\end{rgt}

\begin{orig}
We now extend to include \(J_2\) common close observations after
the treatment period. The total dimension is \(J = J_0 + J_1 + J_2\).

\end{orig}

During the common close phase (time points \(J_1+1, \ldots,
J_1+J_2\)), all participants are off-treatment. The change scores
are:
\begin{align}
  c_{i,J_1+s} &= [\beta(J_1+s) + \gamma g_i \cdot J_1 + b_i(J_1+s)]
    \,t + (w_{i,J_1+s} - w_{i0})
    & s &= 1,\ldots,J_2
\end{align}

\begin{bullets}
\textbf{In the common close the treatment effect persists only as an accumulated offset, fixing the design-matrix entry at the treatment-period value.}
\begin{itemize}\setlength{\itemsep}{2pt}\setlength{\parskip}{0pt}
\item $\gamma$ enters only through the effect accumulated during the treatment period.
\item It does not contribute to the slope during the common close.
\item The $\gamma$ column therefore equals $J_1 t$ for all common-close observations in the treatment group.
\end{itemize}
\end{bullets}

\begin{rgt}
To be completed by rgt.
\end{rgt}

\begin{orig}
The treatment effect \(\gamma\) enters the common close change
scores only through the accumulated effect during the
treatment period (\(\gamma g_i \cdot J_1 \cdot t\)), not through
the slope during the common close. This means the \(\gamma\)
column in the design matrix has the value \(J_1 t\) for all common
close observations in the treatment group.

\end{orig}

The fixed effects design matrix becomes:
\[
  X_i^{(\text{trt})} = t\begin{pmatrix}
    -J_0   & 0 & 0 \\
    \vdots & \vdots & \vdots \\
    -1   & 0 & 0 \\
    0    & 1 & 1 \\
    \vdots & \vdots & \vdots \\
    0    & J_1 & J_1 \\
    0    & J_1+1 & J_1 \\
    \vdots & \vdots & \vdots \\
    0    & J_1+J_2 & J_1
  \end{pmatrix}
\]

\begin{bullets}
\textbf{The slope and treatment columns of the design matrix encode the post-randomization times and the capped treatment exposure respectively.}
\begin{itemize}\setlength{\itemsep}{2pt}\setlength{\parskip}{0pt}
\item The $\beta$ column holds the post-randomization times of the placebo slope across all later visits.
\item The $\gamma$ column rises through the treatment period then stays at $J_1$ during the common close.
\item The $\gamma$ column is identically zero for placebo participants.
\end{itemize}
\end{bullets}

\begin{rgt}
To be completed by rgt.
\end{rgt}

\begin{orig}
The \(\beta\) column contains the post-randomization time for
the placebo slope: \(0,\ldots,0,1,\ldots,J_1,J_1+1,\ldots,J_1+J_2\).
The \(\gamma\) column is \(0,\ldots,0,1,\ldots,J_1,J_1,\ldots,J_1\)
for treatment participants and all zeros for placebo.

\end{orig}

The random effects design remains \(Z_i = t \cdot \mathbf{d}\)
where \(\mathbf{d} = (-J_0,\ldots,-1,1,\ldots,J_1,J_1+1,\ldots,J_1+J_2)'\).

The residual covariance is still
\(R = a(\mathbf{I}_J + \mathbf{1}_J\mathbf{1}_J')\) with
\(J = J_0 + J_1 + J_2\).

\begin{bullets}
\textbf{The common-close extension reuses the same Woodbury computation with only three quantities altered.}
\begin{itemize}\setlength{\itemsep}{2pt}\setlength{\parskip}{0pt}
\item The dimension $J$ increases to include the common-close visits.
\item The $X$ matrix changes only in its $\gamma$ column.
\item The elements of the random-effects vector $Z$ are extended accordingly.
\end{itemize}
\end{bullets}

\begin{rgt}
To be completed by rgt.
\end{rgt}

\begin{orig}
The Woodbury computation proceeds identically, with the only
changes being the dimension \(J\), the structure of the \(X\)
matrix (specifically the \(\gamma\) column), and the elements
of \(Z\).

\end{orig}

We implement the general computation in R below.

\section{Power and Sample Size Formulas}\label{power-and-sample-size-formulas}

For a two-sided test at significance level \(\alpha\) with
power \(1-\beta_{\text{pow}}\), the required number of
participants per group is
\begin{equation}
  m = \frac{(z_{\alpha/2} + z_{1-\beta_{\text{pow}}})^2
    \cdot \mathop{\mathrm{Var}}_1(\hat{\gamma})}{\Delta^2}
  \label{eq:sample_size}
\end{equation}
where \(\mathop{\mathrm{Var}}_1(\hat{\gamma})\) denotes the per-participant-pair
variance from the formulas derived above, and \(\Delta\) is the
minimum clinically important treatment effect on the rate of
change.

Equivalently, the power for a given sample size is
\begin{equation}
  1 - \beta_{\text{pow}} = \Phi\!\left(
    \frac{\Delta\sqrt{m}}
      {\sqrt{\mathop{\mathrm{Var}}_1(\hat{\gamma})}}
    - z_{\alpha/2}
  \right)
  \label{eq:power}
\end{equation}

The relative efficiency of a design with \((J_0, J_1, J_2)\)
observations compared to a standard design with \((0, J_1, 0)\)
is
\begin{equation}
  \text{RE}(J_0, J_1, J_2)
  = \frac{\mathop{\mathrm{Var}}_1(\hat{\gamma})\big|_{0,J_1,0}}
         {\mathop{\mathrm{Var}}_1(\hat{\gamma})\big|_{J_0,J_1,J_2}}
  \label{eq:RE}
\end{equation}

Values greater than 1 indicate that the run-in/common close
design requires fewer participants.

\begin{bullets}
\textbf{A single run-in observation can be counterproductive because the extra slope parameter it introduces may offset the information it contributes.}
\begin{itemize}\setlength{\itemsep}{2pt}\setlength{\parskip}{0pt}
\item Adding one run-in observation introduces the parameter $\delta$, enlarging the fixed-effect dimension.
\item This added degree of freedom can make the one-run-in variance exceed the no-run-in variance for some parameters.
\item Designers should therefore plan either zero or at least two run-in observations, never exactly one.
\end{itemize}
\end{bullets}

\begin{rgt}
To be completed by rgt.
\end{rgt}

\begin{orig}
A practical caveat: adding a single run-in observation
(\(J_0=1\)) introduces the run-in slope parameter \(\delta\)
into the model, increasing the dimension of \(\beta\) from
2 to 3. This additional degree of freedom can offset the
information gained from the run-in observation, so that
\(\mathop{\mathrm{Var}}(\hat\gamma)|_{J_0=1}\) may exceed
\(\mathop{\mathrm{Var}}(\hat\gamma)|_{J_0=0}\) for some parameter
configurations. This effect, noted in a different context
by\textsuperscript{\citeproc{ref-frison1997}{15}}, implies that
trial designers should plan either zero or at least two
run-in observations, not one.

\end{orig}

\section{Numerical Examples}\label{numerical-examples}

\subsection{Verification of closed-form results}\label{verification-of-closed-form-results}

\begin{bullets}
\textbf{The closed-form expressions are validated against the matrix computation in two representative design configurations.}
\begin{itemize}\setlength{\itemsep}{2pt}\setlength{\parskip}{0pt}
\item The verification confirms agreement between the analytic and matrix-based variances.
\item The Frost standard design uses the two-parameter model with two post-baseline visits.
\item The two-run-in case exercises the three-parameter model.
\end{itemize}
\end{bullets}

\begin{rgt}
To be completed by rgt.
\end{rgt}

\begin{orig}
We verify that the closed-form expressions agree with
the matrix computation. The Frost et al.~(2008) result
corresponds to the standard design with \(J_0=0\), \(J_1=2\)
(two post-baseline visits at times \(t\) and \(2t\)), which
under the two-parameter model yields
\(\mathop{\mathrm{Var}}(\hat{\gamma}) = (\sigma^2 + 2\sigma_b^2 t^2)/t^2\).
The \(J_0=2\), \(J_1=1\) case uses the three-parameter model.

\end{orig}

\begin{Shaded}
\begin{Highlighting}[]
\NormalTok{a }\OtherTok{\textless{}{-}} \FloatTok{1.0}
\NormalTok{b }\OtherTok{\textless{}{-}} \FloatTok{0.5}
\NormalTok{tt }\OtherTok{\textless{}{-}} \FloatTok{1.0}

\NormalTok{v\_frost\_matrix }\OtherTok{\textless{}{-}} \FunctionTok{var\_gamma\_matrix}\NormalTok{(}\DecValTok{0}\NormalTok{, }\DecValTok{2}\NormalTok{, }\DecValTok{0}\NormalTok{, a, b, tt)}
\NormalTok{v\_frost\_closed }\OtherTok{\textless{}{-}} \FunctionTok{var\_gamma\_frost}\NormalTok{(a, b, tt)}

\NormalTok{v\_r2\_matrix }\OtherTok{\textless{}{-}} \FunctionTok{var\_gamma\_matrix}\NormalTok{(}\DecValTok{2}\NormalTok{, }\DecValTok{1}\NormalTok{, }\DecValTok{0}\NormalTok{, a, b, tt)}
\NormalTok{v\_r2\_closed }\OtherTok{\textless{}{-}} \FunctionTok{var\_gamma\_r2}\NormalTok{(a, b, tt)}

\FunctionTok{cat}\NormalTok{(}\StringTok{\textquotesingle{}Frost (r=0, k=2) case:}\SpecialCharTok{\textbackslash{}n}\StringTok{\textquotesingle{}}\NormalTok{)}
\end{Highlighting}
\end{Shaded}

\begin{verbatim}
## Frost (r=0, k=2) case:
\end{verbatim}

\begin{Shaded}
\begin{Highlighting}[]
\FunctionTok{cat}\NormalTok{(}\StringTok{\textquotesingle{}  Matrix:\textquotesingle{}}\NormalTok{, }\FunctionTok{round}\NormalTok{(v\_frost\_matrix, }\DecValTok{6}\NormalTok{), }\StringTok{\textquotesingle{}}\SpecialCharTok{\textbackslash{}n}\StringTok{\textquotesingle{}}\NormalTok{)}
\end{Highlighting}
\end{Shaded}

\begin{verbatim}
##   Matrix: 2
\end{verbatim}

\begin{Shaded}
\begin{Highlighting}[]
\FunctionTok{cat}\NormalTok{(}\StringTok{\textquotesingle{}  Closed:\textquotesingle{}}\NormalTok{, }\FunctionTok{round}\NormalTok{(v\_frost\_closed, }\DecValTok{6}\NormalTok{), }\StringTok{\textquotesingle{}}\SpecialCharTok{\textbackslash{}n}\StringTok{\textquotesingle{}}\NormalTok{)}
\end{Highlighting}
\end{Shaded}

\begin{verbatim}
##   Closed: 2
\end{verbatim}

\begin{Shaded}
\begin{Highlighting}[]
\FunctionTok{cat}\NormalTok{(}\StringTok{\textquotesingle{}  Match:\textquotesingle{}}\NormalTok{,}
    \FunctionTok{all.equal}\NormalTok{(v\_frost\_matrix, v\_frost\_closed), }\StringTok{\textquotesingle{}}\SpecialCharTok{\textbackslash{}n\textbackslash{}n}\StringTok{\textquotesingle{}}\NormalTok{)}
\end{Highlighting}
\end{Shaded}

\begin{verbatim}
##   Match: TRUE
\end{verbatim}

\begin{Shaded}
\begin{Highlighting}[]
\FunctionTok{cat}\NormalTok{(}\StringTok{\textquotesingle{}r=2, k=1 case:}\SpecialCharTok{\textbackslash{}n}\StringTok{\textquotesingle{}}\NormalTok{)}
\end{Highlighting}
\end{Shaded}

\begin{verbatim}
## r=2, k=1 case:
\end{verbatim}

\begin{Shaded}
\begin{Highlighting}[]
\FunctionTok{cat}\NormalTok{(}\StringTok{\textquotesingle{}  Matrix:\textquotesingle{}}\NormalTok{, }\FunctionTok{round}\NormalTok{(v\_r2\_matrix, }\DecValTok{6}\NormalTok{), }\StringTok{\textquotesingle{}}\SpecialCharTok{\textbackslash{}n}\StringTok{\textquotesingle{}}\NormalTok{)}
\end{Highlighting}
\end{Shaded}

\begin{verbatim}
##   Matrix: 4.666667
\end{verbatim}

\begin{Shaded}
\begin{Highlighting}[]
\FunctionTok{cat}\NormalTok{(}\StringTok{\textquotesingle{}  Closed:\textquotesingle{}}\NormalTok{, }\FunctionTok{round}\NormalTok{(v\_r2\_closed, }\DecValTok{6}\NormalTok{), }\StringTok{\textquotesingle{}}\SpecialCharTok{\textbackslash{}n}\StringTok{\textquotesingle{}}\NormalTok{)}
\end{Highlighting}
\end{Shaded}

\begin{verbatim}
##   Closed: 4.666667
\end{verbatim}

\begin{Shaded}
\begin{Highlighting}[]
\FunctionTok{cat}\NormalTok{(}\StringTok{\textquotesingle{}  Match:\textquotesingle{}}\NormalTok{,}
    \FunctionTok{all.equal}\NormalTok{(v\_r2\_matrix, v\_r2\_closed), }\StringTok{\textquotesingle{}}\SpecialCharTok{\textbackslash{}n\textbackslash{}n}\StringTok{\textquotesingle{}}\NormalTok{)}
\end{Highlighting}
\end{Shaded}

\begin{verbatim}
##   Match: TRUE
\end{verbatim}

\begin{Shaded}
\begin{Highlighting}[]
\FunctionTok{cat}\NormalTok{(}\StringTok{\textquotesingle{}Variance comparison across designs:}\SpecialCharTok{\textbackslash{}n}\StringTok{\textquotesingle{}}\NormalTok{)}
\end{Highlighting}
\end{Shaded}

\begin{verbatim}
## Variance comparison across designs:
\end{verbatim}

\begin{Shaded}
\begin{Highlighting}[]
\ControlFlowTok{for}\NormalTok{ (rr }\ControlFlowTok{in} \DecValTok{0}\SpecialCharTok{:}\DecValTok{4}\NormalTok{) \{}
\NormalTok{  v }\OtherTok{\textless{}{-}} \FunctionTok{var\_gamma\_matrix}\NormalTok{(rr, }\DecValTok{2}\NormalTok{, }\DecValTok{0}\NormalTok{, a, b, tt)}
  \FunctionTok{cat}\NormalTok{(}\FunctionTok{sprintf}\NormalTok{(}\StringTok{\textquotesingle{}  r=\%d, k=2: Var = \%.4f}\SpecialCharTok{\textbackslash{}n}\StringTok{\textquotesingle{}}\NormalTok{, rr, v))}
\NormalTok{\}}
\end{Highlighting}
\end{Shaded}

\begin{verbatim}
##   r=0, k=2: Var = 2.0000
##   r=1, k=2: Var = 2.0000
##   r=2, k=2: Var = 1.7910
##   r=3, k=2: Var = 1.5000
##   r=4, k=2: Var = 1.2651
\end{verbatim}

\subsection{Sample size as a function of run-in observations}\label{sample-size-as-a-function-of-run-in-observations}

\begin{table}

\caption{\label{tab:table-runin}Required sample size per group ($J_1=2$, $J_2=0$, $\Delta=0.5$, 90\% power, two-sided $\alpha=0.05$) as a function of the number of run-in observations $J_0$, for three levels of between-subject slope variability.}
\centering
\begin{tabular}[t]{rrrr}
\toprule
$J_0$ & Low & Medium & High\\
\midrule
0 & 51 & 127 & 463\\
1 & 44 & 119 & 214\\
2 & 44 & 91 & 113\\
3 & 43 & 70 & 76\\
4 & 42 & 56 & 59\\
\addlinespace
5 & 40 & 48 & 49\\
6 & 38 & 42 & 43\\
\bottomrule
\end{tabular}
\end{table}

\subsection{Sample size as a function of common close observations}\label{sample-size-as-a-function-of-common-close-observations}

\begin{table}

\caption{\label{tab:table-close}Required sample size per group ($J_0=2$, $J_1=2$, $\Delta=0.5$, 90\% power) as a function of common close observations $J_2$.}
\centering
\begin{tabular}[t]{rrrr}
\toprule
$J_2$ & Low & Medium & High\\
\midrule
0 & 44 & 91 & 113\\
1 & 34 & 69 & 80\\
2 & 29 & 48 & 52\\
3 & 25 & 35 & 37\\
4 & 22 & 27 & 28\\
\bottomrule
\end{tabular}
\end{table}

\subsection{Relative efficiency surface}\label{relative-efficiency-surface}

\begin{figure}
\centering
\pandocbounded{\includegraphics[keepaspectratio,alt={(\#fig:fig\_efficiency)Relative efficiency of run-in and common close designs compared to the standard design (J\_0=0, J\_2=0) with the same number of post-randomization observations (J\_1=2). Parameters: \textbackslash sigma\^{}2 = 1, \textbackslash sigma\_b\^{}2 = 1, t = 1.}]{report_files/report_files/figure-latex/fig_efficiency-1.pdf}}
\caption{(\#fig:fig\_efficiency)Relative efficiency of run-in and common close designs compared to the standard design (\(J_0=0\), \(J_2=0\)) with the same number of post-randomization observations (\(J_1=2\)). Parameters: \(\sigma^2 = 1\), \(\sigma_b^2 = 1\), \(t = 1\).}
\end{figure}

\subsection{Application: Alzheimer's disease imaging trial}\label{application-alzheimers-disease-imaging-trial}

\begin{bullets}
\textbf{The Alzheimer's application uses MIRIAD-derived parameters with the correctly combined within-visit residual, which is essential to reproduce Frost's reported sample sizes.}
\begin{itemize}\setlength{\itemsep}{2pt}\setlength{\parskip}{0pt}
\item The slope variance and combined within-visit residual are taken from Frost's MIRIAD estimates.
\item The example assumes a one-year observation interval and a treatment effect of one quarter.
\item Using only the small measurement-error term understates the residual by two orders of magnitude and yields degenerate sample sizes.
\end{itemize}
\end{bullets}

\begin{rgt}
To be completed by rgt.
\end{rgt}

\begin{orig}
Using variance parameters estimated from the MIRIAD dataset\textsuperscript{\citeproc{ref-frost2008}{2}}: \(\sigma_b^2 = 0.9745\)
(between-subject slope variance) and a combined within-visit
residual
\(\sigma^2 = \sigma_u^2 + \sigma_\varepsilon^2
= 0.3338 + 0.0025 = 0.3363\), formed per Eq. \eqref{eq:resid}
from the within-subject between-visit component
(\(\sigma_u^2 = 0.3338\)) and the additional between-scan
measurement error (\(\sigma_\varepsilon^2 = 0.0025\)) of
Frost's Table II. We take observation interval \(t = 1\) year
and treatment effect \(\Delta = 0.25\) (25\% reduction in
atrophy rate relative to placebo rate of approximately 1.0\%
brain volume loss per year). Using the combined residual,
the conventional design reproduces Frost's reported sample
sizes (Table III); using only the \(0.0025\) measurement-error
term, as in earlier drafts, understates the residual by two
orders of magnitude and yields degenerate sample sizes.

\end{orig}

\begin{table}

\caption{\label{tab:ad-example}Sample size per group for AD imaging trial ($J_1=2$ post-randomization visits, $\Delta = 0.25$, 90\% power). Top 15 configurations by sample size.}
\centering
\begin{tabular}[t]{rrrrr}
\toprule
$J_0$ & $J_1$ & $J_2$ & $m$ per group & Total visits\\
\midrule
4 & 2 & 2 & 49 & 9\\
3 & 2 & 2 & 58 & 8\\
4 & 2 & 1 & 60 & 8\\
2 & 2 & 2 & 69 & 7\\
3 & 2 & 1 & 77 & 7\\
\addlinespace
0 & 2 & 2 & 77 & 5\\
1 & 2 & 2 & 77 & 6\\
4 & 2 & 0 & 78 & 7\\
3 & 2 & 0 & 100 & 6\\
2 & 2 & 1 & 104 & 6\\
\addlinespace
1 & 2 & 1 & 138 & 5\\
2 & 2 & 0 & 145 & 5\\
0 & 2 & 1 & 150 & 4\\
1 & 2 & 0 & 247 & 4\\
\bottomrule
\end{tabular}
\end{table}

\begin{verbatim}
## AD imaging trial sample sizes (per group):
\end{verbatim}

\begin{verbatim}
##   Standard (r=0, k=2, f=0): 385
\end{verbatim}

\begin{verbatim}
##   Run-in   (r=2, k=2, f=0): 145  (62% reduction)
\end{verbatim}

\begin{verbatim}
##   Full     (r=2, k=2, f=2): 69  (82% reduction)
\end{verbatim}

\section{Discussion}\label{discussion}

\begin{bullets}
\textbf{The work formalizes and extends Frost's framework with explicit variance formulas for run-in and common-close designs.}
\begin{itemize}\setlength{\itemsep}{2pt}\setlength{\parskip}{0pt}
\item The results provide explicit treatment-effect variance formulas.
\item They cover designs incorporating both run-in and common-close observations.
\item Three substantive points warrant further discussion.
\end{itemize}
\end{bullets}

\begin{rgt}
To be completed by rgt.
\end{rgt}

\begin{orig}
The results presented here formalize and extend the framework
of\textsuperscript{\citeproc{ref-frost2008}{2}} by providing explicit
variance formulas for the treatment effect estimator under
designs that incorporate run-in and common close observations.
Three points merit discussion.

\end{orig}

\begin{bullets}
\textbf{The benefit of run-in observations is governed by the slope-to-noise ratio and is largest for outcomes such as neuroimaging.}
\begin{itemize}\setlength{\itemsep}{2pt}\setlength{\parskip}{0pt}
\item The efficiency gain depends on the ratio of slope variance to residual variance.
\item When that ratio is large, even one or two run-in observations cut sample sizes substantially.
\item For cognitive outcomes with a less favorable ratio, gains are modest and must be weighed against longer study duration.
\end{itemize}
\end{bullets}

\begin{rgt}
To be completed by rgt.
\end{rgt}

\begin{orig}
\textbf{When run-in observations are most beneficial.} The
efficiency gain from run-in observations depends on the ratio
\(\sigma_b^2 t^2 / \sigma^2\). When between-subject slope
variability is large relative to residual variance (as in
neuroimaging outcomes), even one or two run-in observations
yield substantial sample size reductions. For cognitive
outcomes, where this ratio is typically less favorable, the
gains are more modest and the cost of extended study duration
must be weighed against the sample size reduction.

\end{orig}

\begin{bullets}
\textbf{Common-close observations sharpen the treatment-effect estimator at low cost, with diminishing returns beyond the first few visits.}
\begin{itemize}\setlength{\itemsep}{2pt}\setlength{\parskip}{0pt}
\item They inform individual trajectories without keeping participants on treatment.
\item Both groups contribute to slope-variability estimation once treatment has ceased.
\item Marginal benefit diminishes, but the first one or two visits give meaningful reductions when the slope-to-noise ratio is large.
\end{itemize}
\end{bullets}

\begin{rgt}
To be completed by rgt.
\end{rgt}

\begin{orig}
\textbf{Common close as a low-cost efficiency gain.} The common
close phase provides information about individual trajectories
without requiring participants to remain on treatment. Since
treatment has ceased, both groups contribute to estimation of
the underlying slope variability, effectively tightening the
precision of the treatment effect estimator. The marginal
benefit of each additional common close observation diminishes
(as with run-in observations), but the first one or two common
close visits typically provide meaningful sample size
reductions, particularly when \(\sigma_b^2 t^2 / \sigma^2\) is
large.

\end{orig}

\begin{bullets}
\textbf{The derived formulas are consistent with Wang's two-period model while making the three phases explicit and separately optimizable.}
\begin{itemize}\setlength{\itemsep}{2pt}\setlength{\parskip}{0pt}
\item Wang's common-slope scenario corresponds to setting the run-in and placebo slopes equal.
\item Their different-slope scenario corresponds to the general case treated here.
\item The present formulation separates run-in, treatment, and common-close phases, enabling per-phase optimization.
\end{itemize}
\end{bullets}

\begin{rgt}
To be completed by rgt.
\end{rgt}

\begin{orig}
\textbf{Relationship to the two-period model.} The formulas
derived here are consistent with the two-period linear mixed
effects model of\textsuperscript{\citeproc{ref-wang2019}{3}}. Their Scenario 1
(common slope across periods) corresponds to setting
\(\delta = \beta\) in our model, while their Scenario 2
(different slopes) corresponds to the general case. The
advantage of our formulation is the explicit separation of
run-in, treatment, and common close phases, enabling direct
optimization of the number of observations in each phase.

\end{orig}

\begin{bullets}
\textbf{The derivations rest on simplifying assumptions whose relaxation is feasible but typically forfeits closed-form simplicity.}
\begin{itemize}\setlength{\itemsep}{2pt}\setlength{\parskip}{0pt}
\item They assume equally spaced observations, linear trajectories, and a random intercept-plus-slope structure.
\item Extensions to unequal spacing, nonlinear progression, or richer random effects remain within the Woodbury framework but lose simple closed forms.
\item No dropout is assumed, though the Dawson-Lagakos pattern-mixture approach can adjust for anticipated attrition.
\end{itemize}
\end{bullets}

\begin{rgt}
To be completed by rgt.
\end{rgt}

\begin{orig}
\textbf{Limitations.} The derivations assume equally spaced
observations, linear trajectories, and a random intercept
plus random slope covariance structure. Extensions to unequal
spacing, nonlinear progression, and more complex random
effects structures (e.g., random intercept-slope correlation)
are possible within the same Woodbury identity framework but
yield more complex expressions that may not admit simple
closed forms. The formulas also assume no dropout; in
practice, the Dawson-Lagakos pattern-mixture approach\textsuperscript{\citeproc{ref-frost2008}{2}} can be applied to adjust
for anticipated attrition.

\end{orig}

\section{Conflict of interest}\label{conflict-of-interest}

The author declares no conflicts of interest.

\section{Funding}\label{funding}

This work received no specific funding.

\section{Data availability statement}\label{data-availability-statement}

\begin{bullets}
\textbf{All artefacts supporting the paper are openly available in the project repository.}
\begin{itemize}\setlength{\itemsep}{2pt}\setlength{\parskip}{0pt}
\item These include the R package, the four companion manuscripts, and the MATLAB derivation scripts.
\item The shared bibliography is likewise available.
\item Specific paths are enumerated in the Reproducibility section.
\end{itemize}
\end{bullets}

\begin{rgt}
To be completed by rgt.
\end{rgt}

\begin{orig}
The \texttt{runinpower} R package, the four companion manuscripts
of this compendium, the MATLAB symbolic-derivation scripts,
and the shared bibliography are openly available in the
project repository. Specific paths are listed in the
Reproducibility section below.

\end{orig}

\section{Reproducibility}\label{reproducibility}

\begin{bullets}
\textbf{The manuscript build relies on the in-package source together with a small set of testing and rendering packages.}
\begin{itemize}\setlength{\itemsep}{2pt}\setlength{\parskip}{0pt}
\item Rendering used R version 4.4.x.
\item The principal dependencies are the in-package source plus tinytest, rmarkdown, and bookdown.
\item Supporting MATLAB symbolic derivations reside in the analysis scripts directory.
\end{itemize}
\end{bullets}

\begin{rgt}
To be completed by rgt.
\end{rgt}

\begin{orig}
This manuscript was rendered with R 4.4.x. Principal packages
used are the in-package \texttt{runinpower} source (nine source
files: allocation, ancova, averaged, covariance,
covariance-ar1, design, power, replicated, variance),
\texttt{tinytest} (testing harness), and \texttt{rmarkdown} plus
\texttt{bookdown} for the manuscript build. MATLAB symbolic
derivations supporting the Woodbury-based variance formulas
are at \texttt{analysis/scripts/}.

\end{orig}

\begin{bullets}
\textbf{The variance results are deterministic, and any Monte Carlo benchmarking uses a fixed seed for reproducibility.}
\begin{itemize}\setlength{\itemsep}{2pt}\setlength{\parskip}{0pt}
\item Given fixed input parameters the formulas exercise no random-number generation.
\item Monte Carlo benchmarking against the analytic formulas is the only stochastic component.
\item That benchmarking driver fixes the seed at its entry point.
\end{itemize}
\end{bullets}

\begin{rgt}
To be completed by rgt.
\end{rgt}

\begin{orig}
\textbf{Random-number generator.} The variance formulas are
deterministic given fixed input parameters; no RNG is
exercised. Where Monte Carlo benchmarking against the
analytic formulas is reported, the driver sets
\texttt{set.seed(20260418)} at its entry point.

\end{orig}

\begin{bullets}
\textbf{The repository organizes the package source, derivation scripts, companion manuscripts, and bibliography into well-defined locations.}
\begin{itemize}\setlength{\itemsep}{2pt}\setlength{\parskip}{0pt}
\item The R package source and the MATLAB scripts each have a dedicated directory.
\item The three companion manuscripts share a report subtree.
\item This manuscript source and the shared bibliography sit in the run-in power report directory.
\end{itemize}
\end{bullets}

\begin{rgt}
To be completed by rgt.
\end{rgt}

\begin{orig}
\textbf{Data and code availability.} The R package source is at
\texttt{R/}. MATLAB scripts are at \texttt{analysis/scripts/}. The
companion manuscripts are at
\texttt{analysis/report/\{02,03,04\}-*/}. The manuscript source and shared
bibliography (\texttt{../references.bib}) are at
\texttt{analysis/report/01-runin-power/}.

\end{orig}

\section{References}\label{references}

\protect\phantomsection\label{refs}
\begin{CSLReferences}{0}{1}
\bibitem[\citeproctext]{ref-packer2017}
\CSLLeftMargin{1. }%
\CSLRightInline{Packer M. Why has a run-in period been a design element in most landmark clinical trials? {Analysis} of the critical role of run-in periods in drug development. \emph{Journal of Cardiac Failure}. 2017;23(10):697-699. doi:\href{https://doi.org/10.1016/j.cardfail.2017.07.009}{10.1016/j.cardfail.2017.07.009}}

\bibitem[\citeproctext]{ref-frost2008}
\CSLLeftMargin{2. }%
\CSLRightInline{Frost C, Kenward MG, Fox NC. Optimizing the design of clinical trials where the outcome is a rate. {Can} estimating a baseline rate in a run-in period increase efficiency? \emph{Statistics in Medicine}. 2008;27(18):3717-3731. doi:\href{https://doi.org/10.1002/sim.3280}{10.1002/sim.3280}}

\bibitem[\citeproctext]{ref-wang2019}
\CSLLeftMargin{3. }%
\CSLRightInline{Wang G, Aschenbrenner AJ, Li Y, et al. Two-period linear mixed effects models to analyze clinical trials with run-in data when the primary outcome is continuous: {Applications} to {Alzheimer's} disease. \emph{Alzheimer's \& Dementia: Translational Research \& Clinical Interventions}. 2019;5:450-457. doi:\href{https://doi.org/10.1016/j.trci.2019.08.002}{10.1016/j.trci.2019.08.002}}

\bibitem[\citeproctext]{ref-lu2009}
\CSLLeftMargin{4. }%
\CSLRightInline{Lu K, Mehrotra DV, Liu GF. Sample size determination for constrained longitudinal data analysis. \emph{Statistics in Medicine}. 2009;28(4):679-699. doi:\href{https://doi.org/10.1002/sim.3507}{10.1002/sim.3507}}

\bibitem[\citeproctext]{ref-lu2010}
\CSLLeftMargin{5. }%
\CSLRightInline{Lu K. On efficiency of constrained longitudinal data analysis versus longitudinal analysis of covariance. \emph{Biometrics}. 2010;66(3):891-896. doi:\href{https://doi.org/10.1111/j.1541-0420.2009.01322.x}{10.1111/j.1541-0420.2009.01322.x}}

\bibitem[\citeproctext]{ref-hu2021}
\CSLLeftMargin{6. }%
\CSLRightInline{Hu N, Mackey H, Thomas R. Power and sample size for random coefficient regression models in randomized experiments with monotone missing data. \emph{Biometrical Journal}. 2021;63(4):806-824. doi:\href{https://doi.org/10.1002/bimj.202000184}{10.1002/bimj.202000184}}

\bibitem[\citeproctext]{ref-iddi2022}
\CSLLeftMargin{7. }%
\CSLRightInline{Iddi S, Donohue MC. Power and sample size for longitudinal models in {R} -- the longpower package and {Shiny} app. \emph{The R Journal}. 2022;14(1):264-282. \url{https://journal.r-project.org/articles/RJ-2022-022/}}

\bibitem[\citeproctext]{ref-ard2011}
\CSLLeftMargin{8. }%
\CSLRightInline{Ard MC, Edland SD. Power calculations for clinical trials in {Alzheimer's} disease. \emph{Journal of Alzheimer's Disease}. 2011;21:369-377. doi:\href{https://doi.org/10.3233/JAD-2011-0062}{10.3233/JAD-2011-0062}}

\bibitem[\citeproctext]{ref-diggle2002}
\CSLLeftMargin{9. }%
\CSLRightInline{Diggle PJ, Heagerty P, Liang KY, Zeger SL. \emph{Analysis of Longitudinal Data}. 2nd ed. Oxford University Press; 2002.}

\bibitem[\citeproctext]{ref-liuliang1997}
\CSLLeftMargin{10. }%
\CSLRightInline{Liu G, Liang KY. Sample size calculations for studies with correlated observations. \emph{Biometrics}. 1997;53:937-947. doi:\href{https://doi.org/10.2307/2533554}{10.2307/2533554}}

\bibitem[\citeproctext]{ref-sperling2014}
\CSLLeftMargin{11. }%
\CSLRightInline{Sperling RA, Rentz DM, Johnson KA, et al. The {A4} study: Stopping {AD} before symptoms begin? \emph{Science Translational Medicine}. 2014;6(228):228fs13. doi:\href{https://doi.org/10.1126/scitranslmed.3007941}{10.1126/scitranslmed.3007941}}

\bibitem[\citeproctext]{ref-laird1982}
\CSLLeftMargin{12. }%
\CSLRightInline{Laird NM, Ware JH. Random-effects models for longitudinal data. \emph{Biometrics}. 1982;38(4):963-974. doi:\href{https://doi.org/10.2307/2529876}{10.2307/2529876}}

\bibitem[\citeproctext]{ref-wang2022}
\CSLLeftMargin{13. }%
\CSLRightInline{Wang G, Liu L, Li Y, et al. Proportional constrained longitudinal data analysis models for clinical trials in sporadic {Alzheimer's} disease. \emph{Alzheimer's \& Dementia: Translational Research \& Clinical Interventions}. 2022;8(1):e12286. doi:\href{https://doi.org/10.1002/trc2.12286}{10.1002/trc2.12286}}

\bibitem[\citeproctext]{ref-zhao2022}
\CSLLeftMargin{14. }%
\CSLRightInline{Zhao Y, Edland SD. Power formulas for mixed effects models with random slope and intercept comparing rate of change across groups. \emph{The International Journal of Biostatistics}. 2022;18(1):173-182. doi:\href{https://doi.org/10.1515/ijb-2020-0107}{10.1515/ijb-2020-0107}}

\bibitem[\citeproctext]{ref-frison1997}
\CSLLeftMargin{15. }%
\CSLRightInline{Frison L, Pocock SJ. Linearly divergent treatment effects in clinical trials with repeated measures: Efficient analysis using summary statistics. \emph{Statistics in Medicine}. 1997;16(24):2855-2872. doi:\href{https://doi.org/10.1002/(SICI)1097-0258(19971230)16:24\%3C2855::AID-SIM749\%3E3.0.CO;2-2}{10.1002/(SICI)1097-0258(19971230)16:24\textless2855::AID-SIM749\textgreater3.0.CO;2-2}}

\end{CSLReferences}

\vfill

\begin{center}\rule{0.5\linewidth}{0.5pt}\end{center}

\emph{Rendered on 2026-07-01 at 15:37 PDT.}\\
\emph{Source: \texttt{\textasciitilde{}/prj/res/04-runin-power-analysis/runinpower/analysis/report/01-runin-power/report.Rmd}}

\end{document}
